\documentclass[letterpaper, 11pt,  reqno]{amsart}

\usepackage[margin=1.2in,marginparwidth=1.5cm, marginparsep=0.5cm]{geometry}


\usepackage{amsmath, mathrsfs, amssymb,amscd,amsthm,amsxtra, esint}
\usepackage{tikz} 
%\usepackage{latexsym}
%\usepackage[dvips]{graphics,epsfig}
%\usepackage[notcite,notref]{showkeys}
\usepackage{color}

\usepackage[implicit=true]{hyperref}
%\usepackage{showkeys}

\usepackage{cases}%%%

\usepackage{upgreek}

%\usepackage[subnum]{cases}

%\headheight=8pt
%\topmargin=0pt
%\textheight=624pt
%\textwidth=432pt
%\oddsidemargin=18pt
%\evensidemargin=18pt


%\usepackage{hiromacro}
\allowdisplaybreaks[2]

\sloppy

\hfuzz  = 0.5cm %allows mathformula to wiggle a bit

%\includeonly{biblio1}



\setlength{\pdfpagewidth}{8.50in}
\setlength{\pdfpageheight}{11.00in}



%Couleurs
\definecolor{gr}{rgb}   {0.,   0.69,   0.23 }
\definecolor{bl}{rgb}   {0.,   0.5,   1. }
\definecolor{mg}{rgb}   {0.85,  0.,    0.85}
%\definecolor{gy}{rgb}   {0.8,  0.8,   0.8}
\definecolor{yl}{rgb}   {0.8,  0.7,   0.}
\definecolor{or}{rgb}  {0.7,0.2,0.2}
\newcommand{\Bk}{\color{black}}
\newcommand{\Rd}{\color{red}}
\newcommand{\Gr}{\color{gr}}
\newcommand{\Yl}{\color{yl}}
\newcommand{\Bl}{\color{blue}}
\newcommand{\Mg}{\color{mg}}
\newcommand{\Cy}{\color{cyan}}
\newcommand{\Or}{\color{or}}
\newcommand{\dis}{\displaystyle}


\newcommand{\absnabla}{|\nabla|}
\newcommand{\branabla}{\langle\nabla\rangle}
\newcommand{\mathd}{\mathrm{d}}


\usetikzlibrary{shapes.misc}
\usetikzlibrary{shapes.symbols}
\usetikzlibrary{shapes.geometric}


\tikzset{
	ddot/.style={circle,fill=white,draw=black,inner sep=0pt,minimum size=0.8mm},
	>=stealth,
	}


\tikzset{
	ddot2/.style={circle,fill=black,draw=black,inner sep=0pt,minimum size=0.8mm},
	>=stealth,
	}


\newtheorem{theorem}{Theorem} [section]
\newtheorem{maintheorem}{Theorem}
\newtheorem{lemma}[theorem]{Lemma}
\newtheorem{proposition}[theorem]{Proposition}
\newtheorem{remark}[theorem]{Remark}
\newtheorem{example}{Example}
\newtheorem{exercise}{Exercise}
\newtheorem{definition}[theorem]{Definition}
\newtheorem{corollary}[theorem]{Corollary}
\newtheorem{notation}{Notation}[section]
\newtheorem{claim}[theorem]{Claim}
\newtheorem{hypothesis}[theorem]{Hypothesis}
\newtheorem{observation}[theorem]{Observation}

\newtheorem{oldtheorem}{Theorem}

\renewcommand*{\theoldtheorem}{\Alph{oldtheorem}}


\newtheorem{problem}{Problem}

\newtheorem{conjecture}{Conjecture}[section]
\newtheorem{scheme}{Scheme}[section]


%Lower/Upper bound appears below /above the integral sign
\DeclareMathOperator*{\intt}{\int}
\DeclareMathOperator*{\iintt}{\iint}
\DeclareMathOperator*{\supp}{supp}
\DeclareMathOperator{\med}{med}
\DeclareMathOperator*{\ess}{ess}
\DeclareMathOperator*{\Area}{Area}
\DeclareMathOperator{\MAX}{MAX}


\DeclareMathOperator{\HS}{HS}
\DeclareMathOperator{\com}{com}

\DeclareMathOperator{\Id}{Id}
\DeclareMathOperator{\sgn}{sgn}

\DeclareMathOperator{\Law}{Law}
\DeclareMathOperator{\Ker}{Ker}


%
%Roman I
\newcommand{\1}{\hspace{0.5mm}\text{I}\hspace{0.5mm}}
%Roman II
\newcommand{\II}{\text{I \hspace{-2.8mm} I} }
%%Roman III
%\newcommand{\III}{\text{I \hspace{-2.9mm} I \hspace{-2.9mm} I}}
%
%%Roman IV
%\newcommand{\IV}{\text{I \hspace{-2.9mm} V}}

\newcommand{\I}{\mathcal{I}}

\newcommand{\noi}{\noindent}
\newcommand{\Z}{\mathbb{Z}}
\newcommand{\R}{\mathbb{R}}
\newcommand{\C}{\mathbb{C}}
\newcommand{\T}{\mathbb{T}}
\newcommand{\bul}{\bullet}

\newcommand{\XX}{\mathbf{X}}

\newcommand{\B}{\mathcal{B}}

\let\Re=\undefined\DeclareMathOperator*{\Re}{Re}
\let\Im=\undefined\DeclareMathOperator*{\Im}{Im}

\let\P= \undefined
\newcommand{\P}{\mathbf{P}}

\newcommand{\PP}{\mathbb{P}}

\newcommand{\Q}{\mathbf{Q}}


\newcommand{\E}{\mathbb{E}}
\newcommand{\EE}{\mathcal{E}}


\renewcommand{\H}{\mathcal{H}}

\newcommand{\CC}{\mathcal{C}}


\renewcommand{\L}{\mathcal{L}}


\newcommand{\RR}{\mathcal{R}}



\newcommand{\K}{\mathbf{K}}

\newcommand{\F}{\mathcal{F}}

\newcommand{\Nf}{\mathfrak{N}}
\newcommand{\Mf}{\mathfrak{M}}
\newcommand{\Rf}{\mathfrak{R}}

\newcommand{\NB}{\mathbb{N}}
\newcommand{\low}{\textup{low}}
\newcommand{\high}{\textup{high}}


\def\norm#1{\|#1\|}

\newcommand{\al}{\alpha}
\newcommand{\be}{\beta}
\newcommand{\dl}{\delta}
\newcommand{\updl}{\updelta}

\newcommand{\nb}{\nabla}



\newcommand{\Dl}{\Delta}
\newcommand{\eps}{\varepsilon}
\newcommand{\kk}{\kappa}
\newcommand{\g}{\gamma}
\newcommand{\G}{\Gamma}
\newcommand{\ld}{\lambda}
\newcommand{\Ld}{\Lambda}
\newcommand{\s}{\sigma}
\newcommand{\Si}{\Sigma}
\newcommand{\ft}{\widehat}
\newcommand{\Ft}{{\mathcal{F}}}
\newcommand{\wt}{\widetilde}
\newcommand{\cj}{\overline}
\newcommand{\dx}{\partial_x}
\newcommand{\dt}{\partial_t}
\newcommand{\dd}{\partial}
\newcommand{\invft}[1]{\overset{\vee}{#1}}
\newcommand{\lrarrow}{\leftrightarrow}
\newcommand{\embeds}{\hookrightarrow}
\newcommand{\LRA}{\Longrightarrow}
\newcommand{\LLA}{\Longleftarrow}
\newcommand{\LLRA}{\Longleftrightarrow}



\newcommand{\ta}{\theta}

\renewcommand{\l}{\ell}
\renewcommand{\o}{\omega}
\renewcommand{\O}{\Omega}


\newcommand{\bi}{\bfseries\itshape} 

\newcommand{\Gdl}{\mathcal{G}_{\dl} }
\newcommand{\Tdl}{\mathcal{T}_{\dl} }
\newcommand{\Qdl}{\mathcal{Q}_{\dl} }

\newcommand{\Gd}{\wt{\mathcal{G}}_\dl}

\newcommand{\les}{\lesssim}
\newcommand{\ges}{\gtrsim}



\newcommand{\wto}{\rightharpoonup}

%Japanese Bracket
\newcommand{\jb}[1]
{\langle #1 \rangle}


%\renewcommand{\b}{\beta}
%\newcommand{\ind}{\mathbb 1}
\newcommand{\ind}{\mathbf 1}

\renewcommand{\S}{\mathcal{S}}



\newcommand{\gf}{\mathfrak{g}}
\newcommand{\af}{\mathfrak{a}}


\newcommand{\pa}{\partial}
\newcommand{\pP}{\mathbf{P}}
\newcommand{\M}{\mathcal{M}}
\def\e{\eps}

\newcommand{\N}{\mathbb{N}}

%%%%%%%%%%%%%%%%%%%%%%%%
%%%%%%%%%%%%%%%%%%%%%%%%
%%%%%%%%%%%%%%%%%%%%%%%%



\newcommand{\NN}{\mathcal{N}}
\newcommand{\cL}{\mathcal{L}}
\newcommand{\cC}{\mathcal{C}}
\newcommand{\cS}{\mathcal{S}}
\newcommand{\cX}{\mathcal{X}}
\newcommand{\cY}{\mathcal{Y}}
\newcommand{\cZ}{\mathcal{Z}}

\newcommand{\cW}{\mathscr{W}}
\newcommand{\cH}{\mathscr{H}}
\newcommand{\cD}{\mathscr{D}}


%%%%%%%%%%%%%%%%%%%%%%%%
%%%%%%%%%%%%%%%%%%%%%%%%
%New Macros

\newcommand{\W}{\mathcal{W}}

\newcommand{\Lip}{\mathrm{Lip}}
\newcommand{\uw}{U^w}
\newcommand{\uu}{\mathbf{u}}
\newcommand{\vv}{\mathbf{v}}
\newcommand{\z}{\zeta}
\newcommand{\cge}{C^{\g-\varepsilon}}
\newcommand{\tv}{\tilde{v}}
\newcommand{\Ta}{\Theta}
\newcommand{\tI}{\tilde{\mathcal{I}}}
\newcommand{\unn}{\mathbf{u}^{n,N}}
\newcommand{\sub}{\substack}
\newcommand{\tsb}{\theta_{s,\beta}}


\newcommand{\BO}{\text{\rm BO} }
\newcommand{\KDV}{\text{\rm KdV} }
\newcommand{\MKDV}{\text{\rm mKdV} }

\newcommand{\NLS}{\text{\rm NLS} }
\newcommand{\dNLS}{\text{\rm dNLS} }

\newcommand{\too}{\longrightarrow}


\newtheorem*{ackno}{Acknowledgements}



\numberwithin{equation}{section}
\numberwithin{theorem}{section}

\newcommand{\notemax}[1]{\marginpar{\color{red} $\Leftarrow\Leftarrow\Leftarrow$}{\color{red}MAX: #1}}
\newcommand{\noteherbert}[1]{\marginpar{\color{red} $\Leftarrow\Leftarrow\Leftarrow$}{\color{purple}Herbert: #1}}
\newcommand{\hiro}[1]{\marginpar{\color{red} $\Leftarrow\Leftarrow\Leftarrow$}{\smallskip \noi\color{blue}HIRO: #1}\smallskip}
\newcommand{\gp}[1]{\marginpar{\color{red} $\Leftarrow\Leftarrow\Leftarrow$}{\smallskip \noi\color{purple}GP: #1}\smallskip}
\newcommand{\jw}[1]{\marginpar{\color{red} $\Leftarrow\Leftarrow\Leftarrow$}{\smallskip \noi\color{violet}JW: #1}\smallskip}




\makeatletter
\@namedef{subjclassname@2020}{\textup{2020} Mathematics Subject Classification}
\makeatother



\begin{document}
\baselineskip = 14pt



\title[Nonlinear PDEs with modulated dispersion IV]
{Nonlinear PDEs with modulated dispersion IV:\\
normal form approach and unconditional uniqueness}


\author[M.~Gubinelli, G.~Li, J.~Li, and T.~Oh]
{Massimiliano Gubinelli, Guopeng Li, Jiawei Li,  and Tadahiro Oh}


%




\address{Massimiliano Gubinelli\\
Mathematical Institute\\ University of Oxford\\ United Kingdom}

\email{gubinelli@maths.ox.ac.uk}




%


\address{Guopeng Li, 
School of Mathematics and Statistics, Beijing Institute of Technology, Beijing 100081, China}

\email{guopeng.li@bit.edu.cn}



\address{Jiawei Li, School of Mathematics\\
The University of Edinburgh\\
and The Maxwell Institute for the Mathematical Sciences\\
James Clerk Maxwell Building\\
The King's Buildings\\
Peter Guthrie Tait Road\\
Edinburgh\\ 
EH9 3FD\\
 United Kingdom}

 
\email{jiawei.li@ed.ac.uk}


%
\address{
Tadahiro Oh, 
School of Mathematics\\
The University of Edinburgh\\
and The Maxwell Institute for the Mathematical Sciences\\
James Clerk Maxwell Building\\
The King's Buildings\\
Peter Guthrie Tait Road\\
Edinburgh\\
EH9 3FD\\
 United Kingdom\\
and School of Mathematics and Statistics, Beijing Institute of Technology, Beijing 100081, China}


%%
%\address{
%Tadahiro Oh, 
%School of Mathematics and Statistics, Beijing Institute of Technology, Beijing 100081, China\\
%and
%School of Mathematics\\
%The University of Edinburgh\\
%and The Maxwell Institute for the Mathematical Sciences\\
%James Clerk Maxwell Building\\
%The King's Buildings\\
%Peter Guthrie Tait Road\\
%Edinburgh\\
%EH9 3FD\\
% United Kingdom}


\email{hiro.oh@ed.ac.uk}



\subjclass[2020]{35Q53, 35Q35, 35Q55, 
60L20,  65M12, 65M15}



\keywords{modulated dispersion; dispersion management;  Korteweg-de Vries equation; 
Benjamin-Ono equation;
Schr\"odinger equation; 
normal form reduction; unconditional uniqueness; Young integral;
regularization by noise}





\begin{abstract}

We study the modulated  Korteweg-de~Vries equation (KdV)
on the circle
with a time non-homogeneous modulation acting on the linear dispersion term.
By adapting the normal form approach to the modulated setting, 
we prove sharp unconditional uniqueness
of solutions to the modulated KdV in $L^2(\T)$
if a modulation is  sufficiently irregular.
For example, this result implies that if the modulation is given by 
a sample path of a fractional Brownian motion 
with Hurst index $0 < H < \frac 25$, the modulated KdV on the circle is 
unconditionally well-posed in $L^2(\T)$.
Our normal form approach provides the construction of solutions
to the modulated KdV
(and the associated nonlinear Young integral)
{\it without} assuming any positive regularity in time.
As an interesting byproduct of our normal form approach, 
we extend the construction of the 
nonlinear Young integral to a much larger class of functions, 
and obtain an improved Euler approximation scheme 
as compared to the classical sewing lemma approach.
We also establish analogous sharp unconditional uniqueness results
for the modulated Benjamin-Ono equation
and the modulated derivative nonlinear Schr\"odinger equation (NLS)
with a quadratic nonlinearity.
In the appendix, 
we 
prove sharp unconditional uniqueness 
of the cubic modulated NLS on the circle in $H^{\frac 16}(\T)$.






\end{abstract}




%\date{\today}
%%
%
\maketitle
%


\tableofcontents














\newpage

\section{Introduction}\label{SEC:1}

\subsection{Modulated dispersive equations}
\label{SUBSEC:1.1}

%with a time non-homogeneous modulation acting on the linear dispersion term

In the works \cite{CG1, CGLLO1}
 with K.\,Chouk, we  
 considered a {\it modulated} dispersive equation of the following form:
 \begin{align}
\begin{cases}
\dt u +  L u \cdot \dt w=  \NN(u)  \\ 
u|_{t = 0} = u_0,
\end{cases}
\qquad ( t, x) \in \R_+ \times \M,
\label{ME1}
\end{align}


\noi
where $\M  = \R$ or $\T = \R/ (2\pi \Z)$,\footnote{By convention, we endow
$\T$ with the normalized Lebesgue measure $ dx_{_\T} =  (2\pi)^{-1}dx$
such that we do not need to carry factors involving $2\pi$.}
$\NN(u)$ denotes the nonlinearity, 
and 
 $w:\R_+\to\R$ is a continuous function of time, 
called a {\it modulation}, 
 acting on the linear dispersion term $Lu$.
For a specific modulation such as that given by a Brownian motion, 
stochastic calculus can be used to study 
the equation \eqref{ME1};  see \cite{DD, DT, DR2, Ste}.
For a general modulation $w$, however, 
such an approach based on stochastic calculus
is not available.
In \cite{CG1, CGLLO1},  
 with K.\,Chouk, we  
 developed 
  a {\it pathwise} approach 
to study the modulated dispersive equation~\eqref{ME1}
by exploiting {\it irregularity} of the modulation function $w$
(see Definition \ref{DEF:ir})
and  the temporal regularity of a solution~$u$.
In particular, 
in \cite{CGLLO1}, 
we established 
{\it regularization-by-noise}
phenomena for the following  modulated 
 Korteweg-de Vries equation (KdV)
 on both the circle and the real line:
\begin{equation}
\label{kdv1}
\dt u+  \dx^3 u \cdot \dt w =\dx u^2.
\end{equation}

\noi
For example, 
we proved that, given {\it any} $s \in \mathbb R$, the modulated KdV \eqref{kdv1} on $\T$ with a sufficiently irregular modulation is locally well-posed in $H^s(\T)$. 
On the one hand, since the 70's, 
there have been intensive research activities
on regularization by noise 
for stochastic differential equations
and stochastic parabolic partial differential equations;
see
 \cite{
GP, 
KR, 
 FGP, GG0, G22, RT, CGLLO1,
 ABLL} for the  references therein.
See also   survey works \cite{Flan, Gess}.
On the other hand, 
prior to our work~\cite{CGLLO1}, 
% rather few results 
%on 
results on regularization by noise 
%were known
for  dispersive equations
were primarily limited to random initial data\,/\,additive noises of super-critical regularity;
see, for example, 
\cite{BO96, BT08, CO, BT3, BOP2, Poc, OP, 
GKO2, DNY2, DNY3, OOT1, Bring, OOT2, BDNY}.
In~\cite{CGLLO1}, 
we also established 
other forms of regularization by noise
such as 
semilinearization and nonlinear smoothing
(of arbitrary order)
for the modulated KdV \eqref{kdv1}
and other modulated dispersive equations
such as the modulated Benjamin-Ono equation (BO):
\begin{equation}
\label{BO}
\dt u-    \H\dx^2 u \cdot \dt w=\dx u^2,
\end{equation}

\noi
where $\H$ denotes the
Hilbert transform
with the Fourier multiplier $- \ind_{n\ne 0}\cdot i \sgn (n)$.
See~\cite{CGLLO1}
for a further discussion on various examples.
We also mention 
recent works
\cite{Tanaka, Robert1, Robert2, Robert3, Robert4}, 
on pathwise well-posedness of  various modulated dispersive equations;
see \cite[Remark~1.14]{CGLLO1}.

In a recent work \cite{CGLLO2} with A.~Chapouto, we studied
pathwise well-posedness of stochastic modulated dispersive equations
with multiplicative noises
and established a new 
regularization-by-noise phenomenon by exploiting the nonlinear interaction between the unknown and the noise.
For example, 
for 
the stochastic KdV with 
a multiplicative fractional-in-time noise in the Young regime, 
%(with a Hurst parameter $H > \frac 12$), 
%
%
%a multiplicative Young noise, 
we showed that 
irregularity of the modulation induces
smoothing on the stochastic convolution
in a pathwise manner, 
where 
a gain of spatial regularity  becomes (arbitrarily) larger for more irregular modulations.
%see \cite{CGLLO2} for further details.








Since $w$ in \eqref{ME1}
is assumed to be  only continuous, 
the modulated dispersive equation~\eqref{ME1} 
is  formal at this point.
To bypass this problem, 
 we studied 
 the following 
 Duhamel formulation (= mild formulation) of \eqref{ME1}
 in \cite{CG1, CGLLO1}:
\begin{equation}
u(t) = U^w(t) u_0 +  U^w({t}) \int_0^t  U^w({t'})^{-1}\NN( u(t') ) dt', 
\label{mild1}
\end{equation}


\noi
where $U^w(t) = e^{- w(t) L }$ denotes the modulated linear propagator
and we impose $w(0) = 0$ such that  
 $U^w(0) = \Id$.\footnote{The normalization $w(0) = 0$
 is not an additional restriction since only the time derivative $\dt w$
 appears in the modulated equation \eqref{ME1}.}
Define
 the 
 modulated interaction representation
$\uu$ of the unknown $u$ by setting
\begin{align}
\uu(t)=\uw(t)^{-1}u(t), 
\label{int1} 
\end{align}

\noi
for which 
the Duhamel formulation \eqref{mild1}
becomes 
\begin{equation}
\uu(t) = u_0 + \int_0 ^t \uw(t')^{-1}\NN(\uw(t') \uu(t'))dt'.
\label{mild2}
\end{equation}

\noi
Our basic strategy 
in \cite{CGLLO1, CGLLO2}
was 
to first make sense of the integral term in~\eqref{mild2}
as a  nonlinear Young integral
$\I^X(\uu)$ with a given driver $X$
(see, for example, \eqref{K1}
for the modulated KdV)
by applying  the sewing lemma  (Lemma \ref{LEM:sew})
from 
the theory of controlled rough paths~\cite{Gub04}
developed by the first author.
We then studied 
the resulting nonlinear Young differential equation:
\begin{equation}
\uu = u_0 + \I^X(\uu)
\label{YDE0}
\end{equation}

\noi
(possibly with an additive\,/\,multiplicative stochastic forcing)
by a contraction argument. 
We point out that 
 the (sufficient) temporal regularity of
the modulated interaction representation~$\uu$
was essential 
in applying the sewing lemma 
to construct
a nonlinear Young integral.


% via the sewing lemma.
%played a crucial role
%in the construction of a nonlinear Young integral via the sewing lemma.
%

%; see Subsection \ref{SUBSEC:YDE}
%for a further discussion.
%The controlled rough path theory, 
%introduced 
%as an alternative formulation of the rough path theory
%due to Lyons
%\cite{Lyons1, Lyons2}, 
%has been applied and extended further in a variety of contexts; 

\medskip

Now, let us turn to the uniqueness property of  solutions.
Recall 
from  \cite{KATO} 
that  a Cauchy problem is  said to be {\it unconditionally (locally) well-posed}
in $H^s(\M)$
if for every initial condition $u_0 \in H^s(\M)$,
there exist $\tau >0$ and a unique solution $u \in C([0, \tau ];H^s(\M))$
with  $u|_{t = 0} = u_0$.
We refer to such uniqueness in the entire class $C([0, \tau ];H^s(M))$, 
without intersecting with any auxiliary function space, as {\it unconditional uniqueness}.
Unconditional uniqueness is a concept of uniqueness which does not depend
on how solutions are constructed
and thus is of fundamental importance.



Given a driver $X$ of temporal regularity $\frac 12 < \g < 1$, 
a solution $\uu$ to \eqref{YDE0} constructed in \cite{CG1, CGLLO1}
belongs to $\CC^\al([0, \tau] ; H^s(\M))$
for some $0 < \al < 1$ with $\g + \al > 1$
%(see \eqref{Ho2a})
such that the nonlinear Young integral $\I^X(\uu)$
makes sense via the sewing lemma.
As such, 
 the uniqueness, at the level of the modulated interaction representation,  holds 
 {\it conditionally}, i.e.~only  in the class 
 $\CC^\al([0, \tau]; H^s(\M))$.\footnote{Note that the uniqueness
 of solutions to modulated dispersive equations
 constructed 
in \cite{Robert2, Robert3, Robert4}
also hold only conditionally, 
namely only in 
 the  (refined) Fourier restriction norm space $X^s_\tau $.}
 




Our main goal in this paper is to present
an alternative solution theory based
on the {\it normal form method}, 
where we do not require any positive  temporal regularity on
the modulated interaction representation $\uu$.
In particular, we construct a solution 
by working  in $C([0, \tau]; H^s(\M))$
{\it without} utilizing any auxiliary function space, 
and thus establish 
the first unconditional uniqueness results
for  modulated dispersive equations.
From the viewpoint of nonlinear Young integration theory, 
our approach 
presents a completely new construction 
of the nonlinear Young integral $\I^X(\uu)$
in \eqref{YDE0}
 without using any positive temporal regularity, 
but {\it by making use of 
the structure \eqref{YDE0}
of the unknown $\uu$}
(just like a controlled rough path).
In this sense, the novelty of 
our normal form approach
can be viewed,  
at a philosophical level, 
as 
a controlled path approach
to nonlinear Young integration
which 
significantly improves any existing theory
in terms of the temporal regularity.
See Subsection~\ref{SUBSEC:comp}
for a further discussion.
We also point out that our
normal form approach
provides an improvement on  
 the time discretization scheme
as compared to  the sewing lemma approach 
discussed in  \cite[Subsection~2.3]{CG1}
in terms of both the convergence rate and 
the required temporal regularity of a solution;
see
Remark~\ref{REM:num}.
See also Remark \ref{REM:NFx}.
Lastly, while the main body of the paper is devoted
to the modulated KdV-type 
equations with derivative quadratic nonlinearities, 
we also consider 
the modulated cubic nonlinear Schr\"odinger equation (NLS)
on the circle
in Appendix \ref{SEC:A}.

















 














\subsection{Irregularity of the modulation}

We recall a particular notion of irregularity of the modulation $w$, 
introduced in~\cite{CatellierGubinelli, CG1},
which allowed us 
 to establish regularization by noise  in various  ODE and PDE contexts
in \cite{CatellierGubinelli,CG1, CGLLO1, CGLLO2}.
For further study on the notion of irregularity, 
see  \cite{GG, RT}.



\begin{definition}
\label{DEF:ir}
\rm

Let $\rho>0$ and  $0 < \g < 1$.
Given $T > 0$, we say that a function $w\in C([0,T];\R)$ is $(\rho,\g)$-irregular 
on the time interval $[0, T]$  if we have
\begin{align}
\|\Phi^w\|_{  \W^{\rho,\g}_T} 
:= \sup_{a\in \R} \sup_{0\leq r < t\leq T} \langle a \rangle^\rho \frac{|\Phi^w_{t,r}(a)|}{|t-r|^\g} 
< \infty, 
\label{rho1}
\end{align}


\noi
where 
\begin{align}
\Phi^w_{t,r}(a)
=\int_r^t e^{i  a w(t') } d t'.
\label{rho2}
\end{align}


\noi
We say that $w$ is 
$(\rho,\g)$-irregular 
on $\R_+$  if it is 
$(\rho,\g)$-irregular 
on $[0, T]$ for each finite $T > 0$.
%Given an interval $I = [a, b] \subset \R_+$, 
%we say that $w$ is 
%$(\rho,\g)$-irregular 
%on the interval $I$ if $w(\cdot + a)$ is 
%$(\rho,\g)$-irregular 
%on $[0, b-a]$.


\end{definition}









In~\cite{CatellierGubinelli}, 
the authors proved that a fractional Brownian motion is a $(\rho,\g)$-irregular function.
 
 
\begin{oldtheorem}
\label{THM:A}

Let $\{W_t\}_{t\in \R_+}$ be a fractional Brownian motion 
of Hurst index $H\in(0,1)$.
Then, 
for any $\rho < \frac{1}{2H}$,  
there exists $\frac 12 < \g < 1$
such that,  with probability one,  the sample paths of 
$W$ are $(\rho,\g)$-irregular on $\R_+$.

\end{oldtheorem}

Theorem \ref{THM:A}
shows that there exist continuous paths which are $(\rho, \g)$-irregular for arbitrarily large $\rho$. 
In a later work~\cite{GG}, it was shown  that $(\rho,\g)$-irregularity is a generic property of H\"older functions of sufficiently low regularity; see \cite[p.\,2418 and Theorem 3.1]{GG}.

\begin{oldtheorem}
\label{THM:B}
Let  $d \ge 1$.
Given  any $0 < \delta < 1$, 
a generic $\dl$-H\"older continuous  function $w \in C^\delta([0,1];\R^d)$ 
 is $(\rho,\g)$-irregular for any $\rho < \frac1{2\delta}$ with  some $\gamma = \gamma(\rho) \in (\frac 12,1)$.
\end{oldtheorem}


Here, ``generic'' is to be understood according to the notion of {\it prevalence}; see~\cite{GG} for details and the  references therein.




Theorems \ref{THM:A} and \ref{THM:B}
do not only provide important examples of $(\rho, \g)$-irregular paths, 
but also imply that 
 our main results
apply 
to the case
where a modulation is given by a fractional Brownian motion
or a generic  $\dl$-H\"older continuous function.
See, for example, \cite[Corollary~1.5]{CGLLO1}.









\subsection{Nonlinear Young differential equations}
\label{SUBSEC:YDE}





In this subsection, 
we briefly go over the nonlinear Young integration approach
in \cite{CGLLO1}
by taking the modulated 
 KdV~\eqref{kdv1} on~$\T$
 as an example.
See also \cite{HK, G23}
for the construction of nonlinear Young integrals in a more general setting.
For simplicity of the presentation, 
we impose
the mean-zero assumption\footnote{For the modulated KdV \eqref{kdv1}, 
if initial data has non-zero mean $\al_0$, then the following Galilean transformation:
$u(t, x)\mapsto u (t, x + 2\al_0t)-\al_0$, 
along with the conservation of the spatial mean (which can be verified
by arguing as in \cite[Section 6]{CGLLO1}), 
transforms
the equation with a non-zero mean into the mean-zero version.
The same comment applies
to the modulated BO \eqref{BO} and the modulated derivative nonlinear
Schr\"odinger equation \eqref{dNLS1}
studied in this paper.
}
 on solutions to
the modulated KdV \eqref{kdv1}
on $\T$.
We also apply the same convention
to  the other modulated dispersive equations 
considered in this paper.\footnote{Except for the cubic modulated nonlinear Schr\"odinger equation
\eqref{xNLS1}
studied in the appendix.}









Given  $\rho > 0 $ and  $\frac 12  < \g < 1$, 
fix a 
 $(\rho,\g)$-irregular function $w$ on an interval $[0, T]$.
In terms of 
the modulated interaction representation $\uu$
defined in~\eqref{int1}, 
the modulated KdV~\eqref{kdv1}
is written as 
\begin{equation}
\dt \uu = \uw(t)^{-1} \dx\big( (\uw(t) \uu)^2\big), 
\label{kdv1a}
\end{equation}

\noi
where 
$\uw (t)=e^{-   w(t)\dx^3}  $
denotes the modulated linear propagator for \eqref{kdv1}.
By writing~\eqref{kdv1a} (with initial data $u_0$) in the integral form, we have
\begin{equation}
\uu(t) = u_0 + \int_0 ^t \uw(t')^{-1}
\dx\big( (\uw(t') \uu)^2\big)dt'.
\label{mild3x}
\end{equation}

\noi
As mentioned in Subsection \ref{SUBSEC:1.1}, the main strategy
in \cite{CGLLO1}
was to construct the integral term 
in~\eqref{mild3x}
as a nonlinear Young integral via the sewing lemma, 
and solve the resulting nonlinear Young differential equation.



Let $X^\KDV$ be 
the 
bilinear driver  associated with the modulated KdV, 
defined by 
 \begin{align}
X^\KDV_{t,r}(f_1,f_2)
=\int_r^t  \uw(t')^{-1}  \partial_x
\big( (\uw(t')  f_1)( \uw(t') f_2 )\big) dt'
\label{K1}
\end{align}


\noi
for $0 \le r < t \le T$, 
where $f_1$ and $f_2$ are  functions  on $\T$.
By taking the Fourier transform, we have
\begin{align}
\F\big(X^\KDV_{t,r} (f_1,f_2)\big) (n)
= in \sum_{ \substack{n_1, n_2 \in \Z^*\\n = n_1+n_2}}
 \Phi^w_{t,r}(\Xi_\KDV (\bar n))
\ft f_1(n_1)  \ft f_2(n_2),
\label{K2}
\end{align}


\noi
where $\Z^* = \Z \setminus \{0\}$, $\Phi^w_{t,r}$ is as in \eqref{rho2}, 
and $\Xi_\KDV (\bar n)$ denotes the {\it resonance function}\footnote{Here, we follow
the terminology in \cite{Tao2}.  
We point out that $\Xi_\KDV (\bar n)$ is also called the modulation function (see \cite{KOY})
but we do not use this latter terminology to avoid
confusion with a modulation function $w$.
}
 for KdV given by 
\begin{align}
\begin{split}
\Xi_\KDV (\bar n) &  = \Xi_\KDV (n,n_1,n_2) 
= - n^3+ n_1^3 +n_2^3\\
&  = - 3n n_1 n_2, 
\end{split}
\label{K3a}
\end{align}


\noi
where the last equality holds under  $n = n_1+ n_2$.
As pointed out in \cite[Section 3]{CGLLO1}, 
local well-posedness of the modulated KdV \eqref{kdv1}
then follows 
once we show that 
the bilinear driver $X^\KDV_{t, r}$ is 
bounded from $H^s(\T) \times
H^s(\T)$ into $H^s(\T)$
with the norm $\sim |t- r|^\g$, 
uniformly in  $0 \le r < t \le T$
(recall that $\g > \frac 12$), 
which in turn follows from 
an elementary computation, 
using \eqref{rho1}, \eqref{K3a}, 
and 
 Lemma \ref{LEM:SUM};
 see
 \cite[Proposition~4.1]{CGLLO1}.
See also Lemma \ref{LEM:nonlin1}.



As a result, we obtained
the following local well-posedness result
for  the modulated KdV~\eqref{kdv1} on $\T$
(\cite[Theorem 1.3\,(i)]{CGLLO1}). 
We also recall the corresponding 
local well-posedness result
for the modulated BO \eqref{BO} on $\T$
(\cite[Theorem 1.11\,(i)]{CGLLO1}). 



\begin{oldtheorem}\label{THM:1}
Given $\rho \ge\frac 12$,  $\frac12< \g < 1$, and $T> 0$, 
let  $w$ be $(\rho,\g)$-irregular on $[0, T]$ in the sense of Definition~\ref{DEF:ir}.

\smallskip

\noi
\textup{(i)}
Suppose that $\rho \ge \frac 12$ and $s\in \R$
satisfy one of the following conditions\textup{:}
\begin{align}
\begin{split}
\textup{(i.a)} &\ \  \tfrac 12 \le \rho \le \tfrac 34 \quad \text{and} \quad 
s > \tfrac 32 - 3 \rho, \\
\textup{(i.b)} &\ \  \rho > \tfrac 34\quad  \text{and} \quad  s \ge - \rho.
\end{split}
\label{nonlin1b}
\end{align}


\noi
Then, the modulated KdV equation \eqref{kdv1}
on  $\T$
is %semilinearly 
locally well-posed in $H^s(\T)$.






\smallskip

\noi
\textup{(ii)}
Suppose that $\rho \ge 1$ and $s\in \R$
satisfy one of the following conditions\textup{:}
\begin{align}
\begin{split}
\textup{(ii.a)} &\ \   \rho =  1 \quad \text{and} \quad s > - \tfrac 12, \\
\textup{(ii.b)} &\ \  \rho >  1 \quad  \text{and} \quad  s \ge - \tfrac 12 \rho.
\end{split}
\label{regBO1}
\end{align}



\noi
Then, the modulated BO equation \eqref{BO}
on  $\T$
is %semilinearly  
locally well-posed in $H^s(\T)$.



\end{oldtheorem}






Theorem \ref{THM:1}
established regularization by noise
in several respects.
First of all, by taking $\rho$ sufficiently large, 
these  modulated equations
become locally well-posed
in $H^s(\T)$ for
any given $s \in \R$, 
even in the regime
where the unmodulated versions
are ill-posed.
See
\cite{KPV96, CKSTT03, 
KT1, M12,  KV19}
and 
\cite{Moli1, Moli2, 
Deng15, 
GKT23, KLV}
for the known well-\,/\,ill-posedness
results
for the (unmodulated) KdV equation:
\begin{equation}
\dt u+  \dx^3 u  =\dx u^2
\label{kdv2}
\end{equation}



\noi
and 
the (unmodulated) BO equation:
\begin{equation}
\label{BO2}
\dt u-    \H\dx^2 u  =\dx u^2,
\end{equation}

\noi
respectively.
The second point is the 
 semilinearization phenomenon.
The unmodulated 
KdV~\eqref{kdv2}
in $H^s(\T)$ for $s <  - \frac 12$
and
the unmodulated BO \eqref{BO2}
in $H^s(\T)$
 for  any $s \in \R$
are known to behave quasilinearly
such that 
the associated solution map 
is not smooth with respect to the $H^s$-topology;
see  \cite{BO97, MST, CCT,  GKT24}.
On the other hand, 
the construction of solutions in 
Theorem \ref{THM:1} 
is 
 based on a contraction argument
 such that the associated solution map
 is smooth (in fact, real analytic), 
thus  semilinearizing 
these equations
by adding  the modulation.


\begin{remark}\rm 
In a recent work \cite{CGLLO2}, we studied
pathwise well-posedness of stochastic modulated dispersive equations
with multiplicative noises
and established a new 
regularization-by-noise phenomenon by exploiting the nonlinear interaction between the unknown and the noise.
Our strategy was 
to combine
the nonlinear Young integral approach in \cite{CGLLO1}
described above
with a suitable adaptation of 
the pathwise construction
of  stochastic convolutions
via 
the random tensor estimate (see \cite{DNY3, OWZ, OW3}) 
and the sewing lemma, 
which is a strategy
developed in \cite{CLO2, COZ}
to study
 stochastic (unmodulated) dispersive equations with multiplicative noises. 
See \cite{CGLLO2} for details.









\end{remark}












\subsection{Unconditional uniqueness
of modulated dispersive PDEs}
\label{SUBSEC:UU}



In this subsection, 
we state our main results
on
unconditional well-posedness
(and also on nonlinear smoothing)
for modulated dispersive equations on the circle. 
Our first theorem is for the modulated KdV~\eqref{kdv1} on~$\T$.

We say that $u$
is a solution to the modulated KdV \eqref{kdv1} on $\T$, 
if it satisfies the following Duhamel formulation:
\begin{equation}
u(t) = U^w(t) u_0 +  U^w({t}) \int_0^t  U^w({t'})^{-1}\dx u^2(t')  dt'.
\label{mild1x}
\end{equation}


\noi
The same comment applies
to the other modulated dispersive equations considered in this paper.





\begin{theorem}\label{THM:UU1}
Given $\rho > \frac 54$,  $\frac12< \g < 1$, and $T> 0$, 
let  $w$ be $(\rho,\g)$-irregular on $[0, T]$ in the sense of Definition~\ref{DEF:ir}.


\smallskip

\noi
\textup{(i) (unconditional well-posedness).} 
Let   $s\ge 0$.
Then, the modulated KdV equation~\eqref{kdv1}
on  $\T$
is unconditionally locally well-posed in $H^s(\T)$.
In particular, 
if, in addition, $w$ is  $(\rho,\g)$-irregular  on $\R_+$, 
then the modulated KdV equation \eqref{kdv1}
on  $\T$
is unconditionally globally well-posed in $H^s(\T)$.

\smallskip

\noi
\textup{(ii) (nonlinear smoothing).} 
Suppose that %$\rho > \frac 54$, 
$s, s_0 \in \R$ with 
$s_0 > s\ge 0$   satisfy
\begin{align*}
%\qquad \text{and}\qquad 
s_0 < s + 2\rho - \frac 52.
\end{align*}



\noi
Given $u_0 \in H^s(\T)$, 
let $u
\in C([0, \tau]; H^s(\T))$ be the  solution 
to the modulated KdV equation~\eqref{kdv1}
on $\T$
with $u|_{t = 0} = u_0$
constructed in Part (i), 
where 
$\tau \in (0,  T]$ denotes the local existence time.
Then, we have 
\begin{align*}
 u - e^{-w(t) \dx^3}u_0 \in C([0, \tau]; H^{s_0}(\T)).
\end{align*}



\end{theorem}

Theorem \ref{THM:UU1}\,(i) establishes the first unconditional uniqueness
result for modulated dispersive PDEs.
As a corollary to Theorem \ref{THM:UU1}\,(i)
with Theorem \ref{THM:A}, we have
the following unconditional well-posedness
result for the modulated KdV \eqref{kdv1} on $\T$, 
where a modulation is given by a fractional Brownian motion.


\begin{corollary}\label{COR:1}
Let $\{W_t\}_{t\in \R_+}$ be a fractional Brownian motion 
with Hurst index $0 < H < \frac 25$.
%Suppose that $s \in \R$ satisfies 
%\begin{align*}
%\textup{(a)} &\ \  \tfrac 23 \le H  <1 \quad \text{and} \quad 
%s > \tfrac 32 - \tfrac{3}{2H} , \\
%\textup{(b)} &\ \  0 < H < \tfrac 23\quad  \text{and} \quad  s > - \tfrac 1{2H}.
%\end{align*}
%
%\noi
Then, 
with probability one, 
the modulated KdV equation \eqref{kdv1}
on  $\T$
with the modulation given by 
a sample path of the
 fractional Brownian motion 
 $\{W_t\}_{t\in \R_+}$
is unconditionally globally well-posed in $H^s(\T)$
for any $s \ge 0$.




\end{corollary}

\begin{remark}\rm
Similar results hold for each of the modulated dispersive equations considered in this paper, with the appropriate choice of exponents, when the modulation function is given by a fractional Brownian motion, or,  thanks to Theorem \ref{THM:B}, when the modulation function is given by a generic (in the sense of prevalence) H\"older function.
\end{remark}








In view of  the quadratic nonlinearity in \eqref{kdv1}, 
the regularity threshold $s \ge 0$ 
in Theorem~\ref{THM:UU1}
is sharp since the $L^2$-regularity is required to make
sense of $\dx u^2(t)$ as a spatial distribution for given $t \ge 0$
(without intersecting with an auxiliary function space).
The same comment on sharpness of the $L^2$-regularity
also applies to the other modulated dispersive equations
%with the quadratic nonlinearity
considered in this paper.
Note that, in order to make sense of the nonlinearity $\dx u^2$
in $H^s(\T)$ for $s < 0$, we need to 
assume
some temporal regularity on $u$
(or rather its modulated interaction representation)
and give a meaning to $\dx u^2$ as a space-time distribution, 
just as in the Fourier restriction norm method due to Bourgain \cite{BO93}.


While the condition on $\rho$, $s$, and $s_0$
in Theorem \ref{THM:UU1}\,(ii) is more restrictive than
the corresponding result
on the nonlinear smoothing in \cite[Theorem 1.15]{CGLLO1}, 
our argument does not rely
on any auxiliary function space such as 
$\cC^\g([0, \tau]; H^s(\T))$
(even for the modulated interaction representation).


We prove Theorem \ref{THM:UU1}
by implementing a normal form reduction
in the current modulated setting.
In a pioneering work \cite{BIT}, 
Babin, Ilyn, and Titi
implemented  a normal form argument
for the (unmodulated) KdV \eqref{kdv2}
and proved  its unconditional well-posedness 
in $L^2(\T)$.\footnote{In \cite{BIT}, 
Babin, Ilyn, and Titi did not make a connection
of their construction of solutions via ``differentiation by parts''
with 
normal form reductions.
Such a connection was made in an explicit manner
in a later work~\cite{GKO}.}
We point out that, while 
two normal form reductions
were needed in the unmodulated setting~\cite{BIT}, 
we  need to perform a normal form reduction
{\it only once} to prove Theorem \ref{THM:UU1}
for the modulated KdV
thanks to  the presence of the $(\rho, \g)$-irregular modulation function $w$,
which can be viewed as a manifestation
of regularization by noise.
In Subsection \ref{SUBSEC:NF}, we go over basic steps of
the normal form reduction.
We present a proof of Theorem \ref{THM:UU1}
in Section \ref{SEC:KDV}.




\medskip


Next, we state our main  result 
for the modulated BO \eqref{BO} on $\T$.




\begin{theorem}\label{THM:UU2}
Given $\rho > \frac 52$,  $\frac12< \g < 1$, and $T> 0$, 
let  $w$ be $(\rho,\g)$-irregular on $[0, T]$ in the sense of Definition~\ref{DEF:ir}.


\smallskip

\noi
\textup{(i) (unconditional well-posedness).} 
Let   $s\ge 0$.
Then, the modulated BO equation~\eqref{BO}
on  $\T$
is unconditionally locally well-posed in $H^s(\T)$.
In particular, 
if, in addition, $w$ is  $(\rho,\g)$-irregular  on $\R_+$, 
then the modulated BO equation \eqref{BO}
on  $\T$
is unconditionally globally well-posed in $H^s(\T)$.

\smallskip

\noi
\textup{(ii) (nonlinear smoothing).} 
Suppose that 
$s, s_0 \in \R$ with 
$s_0 > s\ge 0$   satisfy
\begin{align*}
s_0 \le  \rho - 1
\qquad \text{and}\qquad 
s_0 < s + \rho - \frac 52.
\end{align*}



\noi
Given $u_0 \in H^s(\T)$, 
let $u
\in C([0, \tau]; H^s(\T))$ be the  solution 
to the modulated BO equation~\eqref{BO}
on $\T$
with $u|_{t = 0} = u_0$
constructed in Part (i), 
where 
$\tau \in (0,  T]$ denotes the local existence time.
Then, we have 
\begin{align*}
 u - e^{w(t) \H\dx^2}u_0 \in C([0, \tau]; H^{s_0}(\T)).
\end{align*}



\end{theorem}

%
%As in the case of the modulated KdV (Theorem \ref{THM:UU1}\,(i)), 
%the regularity range $s \ge 0$
%in Theorem \ref{THM:UU2}\,(i)
%for unconditional uniqueness of the modulated BO \eqref{BO}
%is sharp, since the $L^2$-regularity is required
%to make sense of the quadratic nonlinearity



For the unmodulated BO \eqref{BO2}, 
Mo\c{s}incat and Pilod
proved
its unconditional well-posedness
in $H^s(\T)$ for $s >  3 - \sqrt{\frac {33}4} \approx 0.128$;
see also \cite{Ki3}.
Thus, Theorem \ref{THM:UU2}\,(i)
exhibits a strong regularizing effect in the modulated setting, 
regarding  unconditional well-posedness.
Furthermore, as in the modulated KdV case, 
we only need to carry out a normal form reduction
once for the modulated BO \eqref{BO}, whereas, for the unmodulated BO \eqref{BO2}, 
one first needs to apply a gauge transform 
and then carry out normal form reductions twice.
This  can be once again viewed as a manifestation
of regularization by noise.



We present a proof of Theorem \ref{THM:UU2}
in Section \ref{SEC:BO}.

\begin{remark}\rm
Given $0 < \dl < \infty$, 
 consider the 
modulated 
intermediate long wave equation (ILW) on the circle: 
\begin{equation}
\label{ILW1}
\dt u-    \Gdl\dx^2 u \cdot \dt w =\dx u^2, 
\end{equation}

\noi
where the dispersion operator~$\Gdl \dx^2$ is given 
by 
\begin{align*}
\ft{\Gdl \dx^2 f}(n) =
%\ft{\Gdl \dx^2 }(n)
i p_\dl(n)  \ft f(n)
:=
- \ind_{n\ne 0}\cdot i \bigg(n^2 \coth(\dl n)  - \frac{n}{\dl }\bigg) \ft{f}(n), \quad n\in\Z.
%\label{ILW2}
\end{align*}

\noi
See
\cite{S19, Li24, LOZ, CLOP, CFLOP, FLZ, CLO}
for the known well-\,/\,ill-posedness
results
for the (unmodulated) ILW:
\begin{equation*}
\dt u-    \Gdl\dx^2 u  =\dx u^2.
\end{equation*}


Recall from  \cite[(5.22)]{CGLLO1}
that  the resonance function for ILW 
given by 
\begin{align*}
\Xi_\dl (\bar n)
& = \Xi_\dl(n,n_1,n_2) = - p_\dl(n)  +   p_\dl(n_1) +  p_\dl(n_2)
\end{align*}



\noi
satisfies 
\begin{align}
\Xi_\dl(\bar n) =  \Xi_\BO(\bar n) + O(\dl^{-2}), 
\label{IL8}
\end{align}

\noi
where
$\Xi_\BO(\bar n) $
is the resonance function for BO
defined in \eqref{BX4} below.
Then, together with the fact that 
$|\Xi_\BO(\bar n)| \ges \max(|n|, |n_1|, |n_2|)$
under $n = n_1 + n_2$ and $n n_1 n_2 \ne 0$, 
the relation~\eqref{IL8}
states
that high frequency (multilinear) behavior
of (modulated) ILW is essentially given by that of 
(modulated) 
BO.
In~\cite{CGLLO1}, 
we proved well-posedness
of the modulated ILW \eqref{ILW1}
under the same hypothesis as that for the modulated BO
by using this observation.
Similarly, 
 a slight modification of the proof
of Theorem \ref{THM:UU2}
 with \eqref{IL8}
yields the corresponding result
for the modulated ILW \eqref{ILW1}.
We omit details.






\end{remark}









Lastly, we consider the following modulated derivative 
nonlinear Schr\"odinger equation (dNLS) on the circle:
\begin{align}
i \dt u + \dx^2 u \cdot \dt w =  \dx u^2.
\label{dNLS1}
\end{align}









\begin{theorem}\label{THM:UU3}
Given $\rho > \frac 52$,  $\frac12< \g < 1$, and $T> 0$, 
let  $w$ be $(\rho,\g)$-irregular on $[0, T]$ in the sense of Definition~\ref{DEF:ir}.


\smallskip

\noi
\textup{(i) (unconditional well-posedness).} 
Let   $s\ge 0$.
Then, the modulated dNLS equation~\eqref{dNLS1}
on  $\T$
is unconditionally locally well-posed in $H^s(\T)$.
%In particular, 
%if, in addition, $w$ is  $(\rho,\g)$-irregular  on $\R_+$, 
%then the modulated dNLS equation \eqref{dNLS1}
%on  $\T$
%is unconditionally globally well-posed in $H^s(\T)$.

\smallskip

\noi
\textup{(ii) (nonlinear smoothing).} 
Suppose that 
$s, s_0 \in \R$ with 
$s_0 > s\ge 0$   satisfy
\begin{align*}
s_0 \le  \rho - 1
\qquad \text{and}\qquad 
s_0 < s + \rho - \frac 52.
\end{align*}



\noi
Given $u_0 \in H^s(\T)$, 
let $u
\in C([0, \tau]; H^s(\T))$ be the  solution 
to the modulated dNLS equation~\eqref{dNLS1}
on $\T$
with $u|_{t = 0} = u_0$
constructed in Part (i), 
where 
$\tau \in (0,  T]$ denotes the local existence time.
Then, we have 
\begin{align*}
 u - e^{i w(t) \dx^2}u_0 \in C([0, \tau]; H^{s_0}(\T)).
\end{align*}



\end{theorem}








In \cite{CGKO}, Chung, Guo, Kwon, and the fourth author
studied the unmodulated version of~\eqref{dNLS1} on the circle:
\begin{align}
i \dt u + \dx^2 u =  \dx u^2.
\label{dNLS0}
\end{align}

\noi
By implementing an infinite iteration
of normal form reductions introduced in \cite{GKO}
and combining it with a gauge transform
(the Hopf-Cole transformation), 
they proved
small data unconditional global well-posedness in $L^2(\T)$.
We point out that, in the modulated setting, 
 Theorem~\ref{THM:UU3}
follows from one normal form reduction.
Note also that 
 Theorem \ref{THM:UU3}\,(i)
claims only local well-posedness;
see \cite[Appendix A]{CGKO}
for an explicit construction of
a finite-time blowup solution
to the (unmodulated) dNLS \eqref{dNLS0}.
We present 
 a proof of Theorem \ref{THM:UU3}
in  Section~\ref{SEC:NLS}.














\subsection{Normal form approach
for the modulated KdV}
\label{SUBSEC:NF}


In this subsection, we sketch the main idea
of the normal form approach to construct solutions
to modulated dispersive equations
by taking the modulated KdV \eqref{kdv1} as an example.
In the unmodulated setting, 
this normal form approach was introduced  by Babin, Ilyn, and Titi \cite{BIT}
and was further developed in \cite{KO, GKO, OW, Ki2}.
See also Remark \ref{REM:NF1}.














By taking the Fourier transform of \eqref{kdv1a}, we have
\begin{align}
\dt \ft \uu(t, n) 
= in \sum_{ \substack{n_1, n_2 \in \Z^*\\n = n_1+n_2}}
  e^{i \Xi_\KDV (\bar n)w(t)}
\ft \uu(t, n_1)  \ft \uu(t, n_2), \quad n \in \Z^*, 
\label{kdv1b}
\end{align}

\noi
where $\Xi_\KDV (\bar n) $ is 
the resonance function  in \eqref{K3a}.
In the unmodulated setting (i.e.~$w(t) = t$), 
the first step of normal form reductions
in \cite{BIT} was based
on the identity:
\begin{align}
e^{i \Xi_\KDV (\bar n) t} = \frac{\dt e^{i \Xi_\KDV (\bar n) t}}{i \Xi_\KDV (\bar n)}.
\label{NS0}
\end{align}


\noi
In the current modulated setting, 
we replace this identity by the following
simple, but crucial observation:
\begin{align}
\dd_{t_2}  \Phi^w_{t_1, t_2}(\Xi_\KDV (\bar n)) = - e^{i \Xi_\KDV (\bar n)w(t_2)}, 
\label{NS1}
\end{align}


\noi
which follows from  \eqref{rho2} and the fundamental theorem of calculus.
Then, proceeding with a formal computation,\footnote{Namely, 
we do not worry about justifying computations
such as switching the sum
and the integration
in the formal computation presented here.
We justify all the formal steps in Section \ref{SEC:KDV}.}
a normal form reduction (namely, 
integration by parts in time, using
 \eqref{NS1})
yields
\begin{align}
\ft \uu(t, n) - \ft \uu(0, n) 
& = i n  \int_0^t \sum_{\substack{n_1, n_2 \in \Z^*\\n = n_1 + n_2}}
e^{i  \Xi_\KDV (\bar n)w(t')} \ft \uu(t', n_1)\ft \uu(t', n_2) dt'
\notag\\\
& =-  i n  \sum_{\substack{n_1, n_2 \in \Z^*\\n = n_1 + n_2}}
 \int_0^t
 \partial_{t'} \Phi^w_{t, t'}(\Xi_\KDV (\bar n))\ft \uu(t', n_1)\ft \uu(t', n_2) dt'
 \notag\\
 & =
  i n  \sum_{\substack{n_1, n_2 \in \Z^*\\n = n_1 + n_2}}
\Phi^w_{t, 0}(\Xi_\KDV (\bar n))\ft \uu(0, n_1)\ft \uu(0, n_2) 
\label{NF2}
\\
&  \quad + 2 i n  \sum_{\substack{n_1, n_2 \in \Z^*\\n = n_1 + n_2}}
 \int_0^t
 \Phi^w_{t, t'}(\Xi_\KDV (\bar n))
 \dt \ft \uu(t', n_1) \ft \uu(t', n_2) dt' \notag\\
& =: \F_x\big(X^\KDV_{t, 0}(\uu(0), \uu(0))\big)(n)   + A(t, n), 
\notag
\end{align}

\noi
where $X^\KDV_{t, 0}$ is the bilinear driver defined in 
\eqref{K2}.
At the third step, we used 
the fact that $\Phi^w_{t, t} = 0$, 
the product rule for differentiation, and 
symmetry in $n_1$ and $n_2$. %$\ft \uu(t', n_1)$ and $\ft \uu(t', n_2) $.
See Lemma \ref{LEM:nonlin1} below
for the boundedness and nonlinear smoothing property
of the bilinear driver $X^{\KDV}_{t, 0}$.



\begin{remark}\label{REM:NFx}
\rm
We point out 
that, after the integration by parts in \eqref{NF2},  there is 
no boundary term from the right endpoint  since $\Phi^w_{t, t} = 0$.
This  is different from a usual 
application of a
normal form reduction in the unmodulated setting, 
where we also have a non-trivial contribution from 
the right endpoint.
Furthermore, 
in the usual construction of solutions via the normal form approach, 
the boundary terms do not become small even if we take
a short time interval, 
and thus we need to introduce a cutoff parameter
(based on a frequency or modulation size)
to create smallness to carry out a contraction;
see \cite{BIT, KO, GKO}.
In the current modulated setting, however, 
the estimate \eqref{non1c} for the boundary term comes with the factor $t^\g$
and thus there is no need to introduce such a parameter.


Our version of the normal form reduction, 
using the identity \eqref{NS1}, 
also applies to the unmodulated setting, 
giving a simplification of the argument
in the sense described above
(namely, no need to introduce a parameter
to create a small factor on the boundary terms).



\end{remark}
	





Given $0 \le r <  t \le T$, 
define the trilinear operator $\NN^\KDV_{t, r}(f_1, f_2, f_3)$,  acting on functions on~$\T$, 
 by 
\begin{align}
\begin{split}
& \F_x\big(\NN^\KDV_{t, r}(f_1, f_2, f_3)\big)(n)\\
& \quad =   
- 2 n  \sum_{\substack{n_1, n_2, n_3 \in \Z^*\\n = n_{123}}}
 \Phi^w_{t, r}(\Xi_\KDV (n, n_{12}, n_3))
e^{i  \Xi_\KDV (n_{12}, n_1, n_2)w(r)}n_{12} 
\prod_{j = 1}^3\ft f_j(n_j), 
\end{split} 
\label{nox1}
\end{align}


\noi
where   we used the following short-hand notation:
\begin{align}
n_{1\cdots k} = n_1 + \cdots + n_k.
\label{no1}
\end{align}


\noi
See Lemma \ref{LEM:nonlin2} below
for the boundedness and nonlinear smoothing property
of the trilinear operator $\NN_{t, r}^\KDV$.


By substituting~\eqref{kdv1b}
into $A(t, n)$
on the right-hand side of \eqref{NF2},  we have
\begin{align}
A(t, n) = \int_0^t 
\F_x\big(\NN_{t, t'}^\KDV(\uu(t'), \uu(t'), \uu(t'))\big)(n)dt'.
\label{NF3}
\end{align}



\noi
Then, putting 
\eqref{NF2} and \eqref{NF3} together, 
we have
\begin{align*}
\ft \uu(t, n) - \ft \uu(0, n) 
 = \F_x\big(X^\KDV_{t, 0}(\uu(0)  )\big)(n)
 +  \int_0^t 
\F_x\big(\NN^\KDV_{t, t'}(\uu(t'))\big)(n)dt', 
\end{align*}

\noi
where we used the  short-hand notation \eqref{short1}.
Namely, we have
\begin{align}
\begin{split}
\uu(t) 
 =  u_0 + 
 X^\KDV_{t, 0}(u_0 )
+ 
\int_0^t 
\NN^\KDV_{t, t'}(\uu(t'))dt'.
\end{split}
\label{NF4x}
\end{align}


\noi
Hence, 
we arrive at the following {\it normal form equation}
for the 
modulated KdV \eqref{kdv1}:
\begin{align}
\begin{split}
u(t) 
& = \uw(t) u_0 + 
\uw(t) X^\KDV_{t, 0}(u_0 )
%& \quad +  
+ \uw(t)
\int_0^t 
\NN^\KDV_{t, t'}(\uw(t')^{-1}u(t'))dt'.
\end{split}
\label{NF4a}
\end{align}


\noi
The main point of the normal form reduction
is to transform the original
equation~\eqref{kdv1a}
to the normal form equation \eqref{NF4x}
which 
encodes multilinear dispersive smoothing
(coming from the irregularity of $w$)
in an explicit manner
thanks to the presence of 
$\Phi^w_{t, r}(\Xi_\KDV)$
in $X^\KDV_{t, 0}$
and $\NN^\KDV_{t, t'}$
(see 
\eqref{K2}
and
\eqref{nox1}).
Using the boundedness properties
of $X^\KDV_{t, 0}$
and $\NN^\KDV_{t, t'}$
(Lemmas \ref{LEM:nonlin1}
and \ref{LEM:nonlin2}), 
 a standard contraction argument
yields
unconditional local
well-posedness
of \eqref{NF4x} (and thus of \eqref{NF4a})
in $H^s(\T)$, $s \ge 0$, 
under the hypothesis of 
Theorem~\ref{THM:UU1}.
In the proof of 
Theorem \ref{THM:UU1} presented
in Section \ref{SEC:KDV}, 
we justify the formal computation 
in \eqref{NF2}
and show that 
any solution to the modulated KdV \eqref{kdv1}
(namely, satisfying the Duhamel formulation \eqref{mild1x})
in the class $C([0, \tau]; L^2(\T))$
is also a solution to the normal form equation~\eqref{NF4a}, 
which 
in turn allows us
to show that 
a unique solution
to the normal form equation \eqref{NF4a}
is also an unconditionally unique solution to  the original modulated KdV~\eqref{kdv1}.








\begin{remark}\label{REM:NF9}
\rm


(i) 
Consider 
a modulated dispersive PDE \eqref{ME1}, 
where the nonlinearity $\NN(u)$
is $k$-linear.
Let $n_j$ denote the (spatial) frequency 
of the $j$th factor in $\NN(u)$
and set 
\begin{align}
n_{\max} = \max_{j = 1, \dots, k} |n_j|.
\label{max1}
\end{align}

\noi
If the equation is {\it non-resonant}
with a good lower bound on the resonance function
such as  
$|\Xi(\bar n)|\ges n_{\max}^\ta$
for some (small) $\ta > 0$, 
then
a slight modification of the proof of Theorem \ref{THM:UU1}
yields 
unconditional local well-posedness
by one normal form reduction, 
at least when the modulation is sufficiently irregular.



 




\smallskip

\noi
(ii) 
By adapting the normal form approach 
developed in 
\cite{KOY, Ki2, FO}
for (unmodulated) dispersive equations on $\R^d$, 
one can study unconditional well-posedness of modulated 
dispersive equations on $\R^d$.
We, however, do not pursue this issue here.


\end{remark}








\begin{remark}\label{REM:NF1} \rm
The normal form reduction
presented in this subsection
corresponds to the so-called Poincar\'e normal form reduction, 
where there is no resonance.
When there are resonances, it is called
the Poincar\'e-Dulac normal form reduction;
see \cite[Section 1]{GKO}.
See also Appendix 
 \ref{SEC:A}, 
where 
we implement a Poincar\'e-Dulac normal form reduction for the modulated cubic NLS.
In the unmodulated setting, 
the Poincar\'e-Dulac normal form reductions have been 
used in various contexts:
unconditional uniqueness \cite{BIT, KO, GKO, OW, Ki2, Ki3, MP2, OSW}, 
nonlinear smoothing and the nonlinear Talbot effect \cite{ETz, ETz3, ETz2}, 
existence of global attractors~\cite{ETz3}, 
reducibility~\cite{CGKO}, 
improved energy estimates 
in both the deterministic and probabilistic contexts
\cite{TT, GO, OW2, OST, OS}.
We point out that 
the
higher order $I$-method (with correction terms) \cite{CKSTT03, CKSTT4}
also corresponds to the Poincar\'e-Dulac normal form reduction, 
but applied to the equation satisfied by a modified energy.
It may be of interest to investigate 
analogues of these results
in the modulated setting.





\end{remark}



\begin{remark}\rm

Instead of \eqref{NS1}, we may use the following identity:
\begin{align*}
\dt \Phi^w_{t, 0}(\Xi_\KDV (\bar n)) = e^{i \Xi_\KDV (\bar n)w(t)}
\end{align*}


\noi
and proceed with a normal form reduction.
By repeating the computation presented in this subsection, we
then obtain
a slightly different normal form equation:
\begin{align}
\begin{split}
\uu(t) 
 =  u_0 + 
 X^\KDV_{t, 0}(\uu (t) )
- 
\int_0^t 
\NN^\KDV_{t', 0}(\uu(t'))dt'.
\end{split}
\label{NF4y}
\end{align}


\noi
Note that the argument of $\uu (t)$
of $X^\KDV_{t, 0}$ is now evaluated at the right endpoint of the interval
$[0, t]$.
Compare this with \eqref{NF4x}.
While we prove Theorem \ref{THM:UU1} 
using \eqref{NF4a}, 
one may equally use the normal form equation \eqref{NF4y}
to prove Theorem \ref{THM:UU1} .
We, however, point out that
the normal form reduction with \eqref{NS1}
allows us to directly compare our normal form approach 
with the sewing lemma approach to nonlinear Young integrals
developed in \cite{CGLLO1}.
See the next subsection.
%(The normal form reduction with \eqref{NF1}
%corresponds to 
% the sewing lemma approach to nonlinear Young integrals
%but backward in time.)




\end{remark}





\begin{remark}\rm
Consider the following modulated cubic NLS on $\T$:
\begin{align}
i \dt u + \dx^2 u \cdot \dt w= |u|^2u.
\label{NLS2}
\end{align}

\noi
The associated resonance function 
$\Xi_\NLS (\bar n)$  is given by 
\begin{align}
\begin{split}
\Xi_\NLS (\bar n)
& = \Xi_\NLS(n,n_1,n_2, n_3) = 
n^2 - n_1^2 + n_2^2 - n_3^2\\
& = 2(n- n_1) (n - n_3), 
\end{split}
\label{Xi9}
\end{align}


\noi
where the last step holds under $n = n_1 - n_2 + n_3$, 
and thus we see that 
 the resonance function 
$\Xi_\NLS (\bar n)$
 vanishes when $n = n_1$ or $n_3$.
Namely, the modulated cubic NLS \eqref{NLS2} is 
not non-resonant and thus the observation made in 
Remark \ref{REM:NF9}\,(i) 
does not apply.




In \cite{GKO, OW}, 
the fourth author with his co-authors 
studied a normal form approach to 
 the unmodulated cubic NLS  on $\T$:
\begin{align}
i \dt u + \dx^2 u = |u|^2u.
\label{NLS3}
\end{align}

\noi
In order to overcome
the difficulty coming from  the nearly resonant contributions: $|\Xi_\NLS(\bar n)| \ll n_{\max}$, 
where $n_{\max}$ is as in \eqref{max1},
the authors introduced 
 an 
 infinite iteration of normal form reductions, 
which  transformed \eqref{NLS3}
 into the normal form equation, 
involving infinite series of nonlinearities of arbitrarily high degrees.
This allowed them to prove 
sharp unconditional uniqueness 
of  \eqref{NLS3} in $H^s(\T)$, $s \ge \frac 16$, 
and also unconditional uniqueness
of the normal form equation in $L^2(\T)$
and 
in almost critical Fourier-Lebesgue spaces;
see~\cite{OW}
for further details.



In Appendix \ref{SEC:A}, 
we implement a (Poincar\'e-Dulac) normal form reduction for the modulated cubic NLS \eqref{NLS2}, 
using the identity 
\eqref{NS1},
and prove its unconditional uniqueness;
 see Theorems \ref{THM:UU4}. % and \ref{THM:UU5}.
\end{remark}


\subsection{Comparison with 
the sewing lemma approach}
\label{SUBSEC:comp}


In this subsection, we
compare  the normal form approach 
presented in the previous subsection  with the sewing lemma approach 
to nonlinear Young integrals in \cite{CGLLO1}.
For this purpose, we first need to introduce some preliminary notations.


Let $V$ be a Banach space and $T>0$.
For $n\in\N$, we denote 
\begin{align*}
\Delta_{n, T} = 
\big\{ (t_1, \ldots, t_n) \in [0,T]^n: \ t_i > t_j 
\text{ for } i < j\big\}.
\end{align*}
We denote by $C_{n,T}V$ 
the space of continuous functions 
from $\Delta_{n,T}$ to $V$. 
%When $n=1$,
%we may write $C_T V$ for simplicity, 
%and equip this space with the supremum norm: 
%\begin{align*}
%\|f\|_{C_T V} = \|f\|_{L^\infty_T V} = \sup_{0\leq t \leq T} \|f(t)\|_V.
%\end{align*}
%
%
%
%\noi
We define the coboundary operator 
$\updl: C_{n,T} V  \to C_{n+1,T} V$ 
as follows; 
given  $f\in C_{n,T} V$ 
and $(t_1,\ldots, t_{n+1}) \in \Dl_{{n+1},T}$, 
we set
\begin{align*}
(\updl f)_{t_1,\ldots , t_{n+1}} = 
\sum_{k=1}^{n+1} (-1)^{n-k} f_{t_1, \ldots, t_{k-1}, t_{k+1}, \ldots, t_{n+1}}.
\end{align*}



\noi
For example, for $f\in C_T V$
and 
$g\in C_{2,T} V$, 
we have 
\begin{align}
\begin{split}
(\updl f )_{t,r} &= f_t - f_r, \\
(\updl g)_{t_1,t_2,t_3} &= 
g_{t_1,t_3} - g_{t_1,t_2} - g_{t_2,t_3}
\end{split}
\label{dl1}
\end{align}



\noi
for $(t,r)\in\Dl_{2,T}$
and $(t_1,t_2,t_3) \in \Dl_{3,T}$.
As noted in \cite{GT10}, 
the sequence 
\begin{align*}
0 \too \R \too C_{1, T}V \stackrel{\updl}{\too} C_{2, T}V
\stackrel{\updl}{\too} C_{3, T}V
\stackrel{\updl}{\too} \cdots
\end{align*}

\noi
is exact. 
In particular, we have 
 $\updl\circ\updl =0$ and 
if $f \in C_{n,T} V$ with $\updl f =0$, 
then there exists a $g\in C_{n-1,T}V$ 
such that $f= \updl g$; 
see, for example,  \cite[Lemma 2.1]{GT10}. 

Given $0 < \g \le 1$, we denote by $C^\g_T V = C^\g([0, T]; V)$ the space of $\g$-H\"older continuous functions taking values in $V$, endowed  with the seminorm:
\begin{align*}
\|f\|_{C^\g_T V} = \sup_{(t,r)\in \Dl_{2,T}} 
\frac{\|(\updl f)_{t,r}\|_V}{|t-r|^\g}.
\end{align*}

\noi
We also define 
 $\CC^\g_T V = \CC^\g([0, T]; V)$ via the norm:
\begin{align}
\| f \|_{\CC^\g_TV} = \| f \|_{L^\infty_TV} + \|f\|_{C^\g_T V}.
\label{Ho1}
\end{align}


\noi 
We also introduce the spaces 
$C^\g_{n,T}V$, $n=2,3$, 
equipped with the following H\"older-type norms;
 for $g\in C_{2,T}V$ and $h\in C_{3,T}V$, we set
\begin{align}
\begin{split}
\| g\|_{C^\g_{2,T} V} & 
= \sup_{(t,r) \in \Dl_{2,T}} \frac{\|g_{t,r} \|_{V} }{|t-r|^{\g}}, \\
\| h\|_{C^\g_{3,T} V} & 
= \inf_{0<\al<\g} 
\sup_{(t_1,t_2,t_3) \in \Dl_{3,T}} 
\frac{ \|h_{t_1,t_2,t_3}\|_V}
{|t_1-t_2|^{\al} |t_2-t_3|^{\g-\al}}.
\end{split}
\label{Ho2}
\end{align}



\noi
Given a Banach space $V$ and $k \in \N$, 
we  use $\Lip_k(V)$ 
to denote the Banach space of 
locally Lipschitz maps $f:V\to V$
with polynomial growth of order $k$
such that 
\begin{align}
\|f\|_{\Lip_k(V)} = 
\sup_{x,y\in V} 
\frac{\|f(x)-f(y)\|_V}{\|x-y\|_V \big(1+\|x\|_V+\|y\|_V\big)^{k-1}}
<\infty .
\label{Lip1}
\end{align}



\noi
Then, 
from \eqref{Ho2} and \eqref{Lip1}, we have 
\begin{align}
\|f \|_{C^\g_{2, T}\Lip_k(V)}
= \sup_{(t,r) \in \Dl_{2,T}}
\frac 1{|t-r|^{\g}}
\sup_{x,y\in V} 
\frac{\| f_{t, r}(x)- f_{t, r}(y)\|_V}{\|x-y\|_V \big(1+\|x\|_V+\|y\|_V\big)^{k-1}}.
\label{Ho3}
\end{align}



\medskip
 
 Next, we 
review the essential idea
 on the construction of a nonlinear Young integral 
 via the sewing lemma (Lemma \ref{LEM:sew}) presented in \cite{CGLLO1}.
Fix a driver  $X\in C^\g_{2,T}\Lip_k(V)$, where $C^\g_{2,T}\Lip_k(V)$
is as in~\eqref{Ho3},  such that 
\begin{align*}
X_{t,r}(0)=0
%\label{Ja0}
\end{align*} 

\noi
for any $(t,r) \in  \Dl_{2,T}$, 
and 
\begin{align}
\label{Ja1}
(\updl X)_{t_1,t_2,t_3} = 0
\end{align}


\noi
for any $(t_1,t_2,t_3)\in \Dl_{3,T}$.
Given 
$\uu\in \CC^{\al}([0,T];V)$ for some $0 < \al < 1$, 
the nonlinear Young integral $\I^X(\uu)$ of $\uu$
is a unique function 
with $\I^X(\uu)(0) = 0$
whose increment is given by 
\begin{align}
\updl 
\big(\I^X(\uu)\big)_{t, r} 
= X_{t, r}(\uu_\bul).
\label{YQ1}
\end{align}

\noi
Here, $\uu_\bul$ denotes 
the function $\uu$ evaluated at the variable of integration
$\bul$, satisfying $r \le \bul \le  t$.
As it is, the expression on the right-hand side of \eqref{YQ1}
is not well defined.
By replacing $\bul$ by the left endpoint $r$, 
we have 
\begin{align}
X_{t, r}(\uu_\bul)
= X_{t, r}(\uu_r) + R_{t, r}
\label{YQ2}
\end{align}

\noi
for some two-parameter process $R = R^{X, \uu}$.
%Once we find one such error term $R$ with sufficient regularity,
%satisfying \eqref{YQ2}, 
%the sewing lemma allows us to  define  the nonlinear Young integral
% $\I^X(\uu)$ in a unique manner.
% 
Define  $\Theta = \Ta^{X, \uu}$ on $\Dl_{2,T}$
by setting
\begin{align}
\Theta_{t,r} = X_{t,r}(\uu_r), \quad (t,r)\in \Dl_{2,T}.
\label{Ja1a}
\end{align}


\noi
Then, 
by applying the coboundary operator $\updl$
defined in \eqref{dl1}
to \eqref{YQ2}
and noting that the left-hand side of~\eqref{YQ2} vanishes
under $\updl$, 
any error term $R$ in~\eqref{YQ2} (if it exists) satisfies
\begin{align}
(\updl R)_{t_1, t_2, t_3}
= - (\updl  \Theta)_{t_1, t_2, t_3}
=  X_{t_1, t_2}( \uu_{t_2}) - 
  X_{t_1, t_2}( \uu_{t_3})
%= X_{t_1, t_2}\big((\updl \uu)_{t_2, t_3}\big).
\label{YQ3}
\end{align}

\noi
for any $(t_1, t_2, t_3) \in\Dl_{3, T}$, 
where we used 
 \eqref{Ja1} at the second equality.



We now recall 
the  sewing lemma due to the first author \cite{Gub04};
see also
\cite[Proposition~2.3, Corollaries  2.4 and  2.5]{GT10}
and  \cite[Lemma~4.2]{FH20}.
%
We set 
$C^{1+}_{n,T}V = \bigcup_{\g>1}C^\g_{n,T}V$.

\begin{lemma}[sewing lemma]
\label{LEM:sew}


Let $V$ be a Banach space and 
fix $T>0$. 
Then,  there exists a unique linear map \textup{(}called the sewing map\textup{)}
$\Lambda:C^{1+}_{3,T} V \cap 
\Ker \updl|_{C_{3,T}V}
\to C^{1+}_{2,T}V$ such that 

\smallskip
\begin{enumerate}
\item[(i)] 
We have 
$\updl \Lambda h = h$
for each  $h\in C_{3,T}V\cap \Ker  \updl|_{C_{3,T} V}$.



\smallskip

\item[(ii)]
 For each $\z >1$,
the sewing map $\Lambda$ is continuous
from $C^\z _{3,T}V\cap 
\Ker  \updl|_{C_{3,T} V}$ to 
$C^\z _{2,T}V$ such that 
\begin{align*}
\|\Lambda h \|_{C^\z _{2,T}V} 
\le \frac{1}{2^\z - 2}  \| h \|_{C^\z _{3,T}V} 
\end{align*}


\noi
for any $h\in C^\z _{3,T}V$.



\smallskip
\item[(iii)] 
Given any  $g\in C_{2,T}V$ 
with 
$\updl g\in C^\z _{3,T}V$, 
 there exists  unique
$f\in C([0, T];V)$  \textup{(}modulo an additive  constant\textup{)} 
such that 
$\updl f = (\Id - \Lambda \updl)g$. 
In addition, 
we have 
\begin{align*}
(\updl f)_{t,r} = \lim_{|\Pi([r,t])|\to 0} 
\sum_{j=0}^n g_{t_j,t_{j+1}}
\end{align*}


\noi
 for any $(t,r)\in \Dl_{2,T}$,
 where 
 the limit is over any partition
 $\Pi([r,t])$  
 of  $[r,t]$\textup{:} 
\[\Pi ([r,t]) = \{r = t_n < \dots < t_1 <  t_0 = t\}\]
whose mesh size 
$|\Pi([r,t])| = \sup_{j} |t_j-t_{j+1}|$ 
tends to $0$.


\end{enumerate}


\end{lemma}







Since
$X\in C^\g_{2,T}\Lip_k(V)$
and 
$\uu\in \CC^{\al}([0,T];V)$, 
it follows from \eqref{YQ3} with \eqref{Ho3} and \eqref{Ho1} that 
$\updl R 
= -\updl \Ta 
 \in C^{\g + \al}_{3,T}V$, 
 where $\Ta$ is as in \eqref{Ja1a}; see \cite[(3.27)]{CGLLO1}.
Hence, if $\g + \al > 1$, 
then we can apply the sewing lemma (Lemma~\ref{LEM:sew})
to define an error term $R$ by setting
\begin{align}
R = - \Ld \updl  \Theta
\in C^{\g + \al}_{2,T}V.
\label{YQ4} 
\end{align}



\noi
Then, 
putting
 \eqref{YQ1},  
 \eqref{YQ2},  
and \eqref{YQ4} together,   
we define the nonlinear Young integral $\I^X(\uu)$ of $\uu$
(with respect to the nonlinear Young driver $X$)
to be  a unique function 
in $ \CC^\g([0,T]; V)$
with $\I^X(\uu)(0) = 0$
whose increment is given by 
\begin{align}
\updl (\I^X(\uu))
= (\Id - \Ld \updl )\Ta.
\label{YQ5}
\end{align}





\noi
Once \eqref{YQ5} holds, 
the nonlinear Young integral $\I^X(\uu)$ 
is given by 
the unique limit of 
Riemann-Stieltjes type sums:
\begin{align}
\I^X(\uu)(t) 
= \lim_{|\Pi([0,t])|\to 0} 
\sum^n_{j=0} \Theta_{t_j,t_{j+1}}
= \lim_{|\Pi([0,t])|\to 0} 
\sum^n_{j=0} X_{t_j,t_{j+1}}(\uu_{t_{j+1}}), 
\label{YQ6}
\end{align} 
 
 
 
\noi
 where 
the limit is in the sense of Lemma \ref{LEM:sew}\,(iii).






\medskip



We are now ready to compare the normal form approach
with the sewing lemma approach.
Given  $\rho > \frac 54 $ and  $\frac 12  < \g < 1$, 
fix a 
 $(\rho,\g)$-irregular function $w$ on an interval $[0, T]$
 and $s\ge 0$ as in Theorem \ref{THM:UU1}.
Consider the modulated KdV \eqref{kdv1} on $\T$
and write 
\eqref{mild3x}
as
\begin{equation}
\uu = u_0 + \I^{X^\KDV}(\uu), 
\label{YZ1}
\end{equation}

\noi
Here, 
$X^\KDV$
is the bilinear driver defined in \eqref{K1}.
The main goal is to give a precise meaning
to 
the integral term $\I^{X^\KDV}(\uu)$.

\medskip

\noi
$\bul$ {\bf Sewing lemma approach:}
From 
Lemma \ref{LEM:nonlin1}, 
we have
$X\in C^\g_{2,T}\Lip_k(V)$.
Thus, 
if $\uu\in \CC^{\al}([0,T];H^s(\T))$
for some $0 < \al < 1$ such that $\g + \al > 1$, 
then
we can apply the sewing lemma 
(Lemma \ref{LEM:sew})
as described above
to construct
the nonlinear Young integral $\I^{X^\KDV}(\uu)$. 
From \eqref{YQ5} with \eqref{Ja1a}, we have 
\begin{align}
\I^{X^\KDV}(\uu)(t) - \I^{X^\KDV}(\uu)(r)
= X^\KDV_{t, r} (\uu(r)) - (\Ld \updl \Ta)_{t, r}
\label{YZ2}
\end{align}

\noi
for any $0 \le r < t \le T$, 
where $\Ld$ is the sewing map.




\medskip

\noi
$\bul$ {\bf Normal form approach:}
Let $\uu \in C([0, \tau]; H^s(\T))$
be the unique solution to \eqref{YZ1}
constructed in Theorem \ref{THM:UU1}.
Then, by repeating the computations in \eqref{NF2}
over an time interval $[r, t]$, 
we have
\begin{align}
\begin{split}
\I^{X^\KDV}(\uu)(t) - \I^{X^\KDV}(\uu)(r)
& = \uu(t) - \uu(r) \\
&  =  
 X^\KDV_{t, r}(\uu (r) )
+ 
\int_r^t 
\NN^\KDV_{t, t'}(\uu(t'))dt'
\end{split}
\label{YZ3}
\end{align}

\noi
for any $0 \le r < t \le \tau$, 
where $\NN^\KDV_{t, t'}$ is as in \eqref{nox1}.
We point out that in obtaining \eqref{YZ3}, 
we used the fact that $\uu$ is a solution to \eqref{YZ1}.


\medskip

By comparing \eqref{YZ2} and \eqref{YZ3}, 
we obtain 
\begin{align}
  (\Ld \updl \Ta)_{t, r}
 = 
- \int_r^t 
\NN^\KDV_{t, t'}(\uu(t'))dt'
\label{YZ4}
\end{align}




\noi
for any $0 \le r < t \le \tau$, 
where
the left-hand side makes sense
only for  $\uu\in \CC^{\al}([0,\tau];H^s(\T))$
with $\g + \al > 1$, 
whereas 
the right-hand side makes sense
for  $\uu\in C([0,\tau];H^s(\T))$.
Hence, we can view 
the right-hand side
of \eqref{YZ4}
as an extension of 
the definition of 
$\Ld \updl \Ta$ on the left-hand side
to functions of lower temporal regularity $\al \ge 0$, 
provided that $\uu$ is a solution to \eqref{YZ1}.
Namely, 
in the current modulated setting, 
{\it the normal form reduction
with the (controlled) structure~\eqref{YZ1}
extends 
the construction of the nonlinear Young integral
$\I^{X^\KDV}(\uu)$ 
to the much larger class 
$C([0,\tau];H^s(\T))$}, 
providing a concrete expression for $\Ld \updl \Ta$.
Moreover, it follows from Lemma \ref{LEM:nonlin2}
that the right-hand side of \eqref{YZ4}
has temporal regularity $1 + \g$
(and, as a result,  so does the left-hand side of~\eqref{YZ4})
which shows an improvement of the temporal regularity
from $\g + \al$ resulting from the application of the sewing lemma (Lemma \ref{LEM:sew}).
As a consequence, 
we see that  
$\I^{X^\KDV}(\uu)$ defined by the right-hand side of \eqref{YZ3}
is still given by the 
 unique limit of 
Riemann-Stieltjes type sums in \eqref{YQ6}.




\begin{remark}\label{REM:num}\rm


We  note that 
the nonlinear Young integral approach in \cite{CG1, CGLLO1, CGLLO2}
allows for a straightforward Euler approximation scheme
to discretize in time (see, for example,  \cite[Subsection~2.3]{CG1};
see also \cite[Subsection 3.3]{GubinelliKdV}).
Similarly, 
by repeating the computation
in \cite[Subsection~2.3]{CG1}
with  \eqref{YZ3} instead of \eqref{YZ2}, 
we obtain an analogous convergence result 
but with a faster convergence rate 
thanks to the higher temporal regularity $1+ \g$
of the second term on the right-hand side of \eqref{YZ3}
(as compared to 
$ - \Ld \updl \Ta$ in \eqref{YZ2}) as pointed above.
Moreover, this approximation scheme \eqref{YZ3} via the normal form approach
does {\it not} assume any positive temporal regularity on a solution,
which is a significant improvement.
%even from the viewpoint of time discretization.

We point out that ideas
analogous to the normal form approach 
(and the sewing lemma approach) have been implemented
in time discretization schemes; see, for example, 
\cite{HSch, BSch}.
We also mention 
the work 
\cite{HMSch} on convergence of a time-discretization scheme
for the modulated 
nonlinear Schr\"odinger equation.
For example, 
the approximation via the first order exponential integrator 
for the (unmodulated) KdV 
in \cite{HSch}
is given by 
\begin{align}
 \uu(t_j) 
\approx
 \uu(t_{j+1}) +  X^\KDV_{t_j, t_{j+1}}(\uu (t_{j+1}) ), 
\label{YZ5}
\end{align}

\noi
where $X^\KDV_{t, r}$
is as in 
\eqref{K1} with $w(t) = t$.
(Recall our convention $t_j > t_{j+1}$.)
Compare~\eqref{YZ5} with \eqref{YZ2} and \eqref{YZ3}, 
where the second terms on the right-hand sides of 
\eqref{YZ2} and \eqref{YZ3}
are now viewed as the error terms
in the discretization scheme \eqref{YZ5}, giving (local-in-time) convergence rates.
Note that \eqref{YZ2} and \eqref{YZ3}
hold only in the modulated setting.\footnote{In the unmodulated setting, 
by assuming a higher regularity for a  solution
to the (unmodulated) KdV, 
we can make sense of the second term 
on the right-hand side of \eqref{YZ3}
and carry out local-in-time and global-in-time error analysis;
see \cite{CFO}
for further details.}
In the unmodulated setting, 
we need to use the second order expansion
in \cite{Gub04} via the nonlinear controlled path approach
(and \cite{BIT} via the second normal form reduction)
to replace \eqref{YZ2} (and \eqref{YZ3}, respectively)
which still provides a convergence result.





\end{remark}



\section{Notations and preliminary tools}
\label{SEC:2}


Let $A\les B$ denote an estimate of the form $A\leq CB$ for some constant $C>0$. We write $A\sim B$ if $A\les B$ and $B\les A$, while $A\ll B$ denotes $A\leq c B$ for some small constant $c> 0$. 
%We may write  $\les_{\al}$ and $\sim_{\al}$ to 
%emphasize the dependence on an external parameter $\al$.
We use $C>0$ to denote various constants, which may vary line by line.



In expressing the dependence of a function $u$
on the time variable, we often use the short-hand notation
$u_t = u(t)$,  which is standard in probability theory and stochastic analysis.





We  use $\ft f$ and $\F_x(f)$
to denote
the  Fourier transform
of a function $f$ on $\T$, 
defined by 
\begin{align*}
\ft  f(n)=\int_{\T}f(x)e^{-i n x} d x_\T
= \frac 1 {2\pi}\int_{\T}f(x)e^{-i n x} d x.
\end{align*}

\noi
We impose the mean-zero assumption
on  the equations we consider in this paper
and thus we set  $\Z^* = \Z \setminus \{0\}$.
Given $s \in \R$, 
we define the  Sobolev space $H^{s}(\T)$ 
via the norm: 
\begin{align*}
\|  f  \|^2_{H^{s}(\T)}
& =
\sum_{n \in \Z} \jb{n}^{2s} |\ft f(n)|^2.
  \end{align*}


We often use short-hand notations such as
$C_T H^s_x  = C\big([0, T]; H^s(\M))$, 
 when there is no ambiguity.





Given $k \in \N$
and Hilbert spaces $H_1, H_2$,  
we use 
\begin{align*}
\cL_k(H_1; H_2)
\end{align*}

\noi
 to denote 
the Banach space of bounded $k$-linear operators 
from  $H_1^{\otimes k}$ into $H_2$.
Given a multilinear operator $S$, 
we use the following short-hand notation:
\begin{align}
S(f) = S(f, \dots, f).
\label{short1}
\end{align}










\medskip







We recall the following basic lemma on a  discrete convolution
which follows
from an elementary  computation.
See, for example,  
 \cite[Lemma 4.2]{GTV}. % and \cite[Lemma 4.1]{MWX}.





\begin{lemma}\label{LEM:SUM}
Let  $\al, \be \in \R$ satisfy
\begin{align*}
 \al \ge \be \ge 0 \qquad \text{and}\qquad  \quad \al+ \be > 1.
\end{align*}


\noi
Then, we have
\begin{align*}
 \sum_{n = n_1 + n_2} \frac{1}{\jb{n_1}^\al \jb{n_2}^\be}
& \les \frac 1{\jb{n}^{ \be - \ld}}
\end{align*}

\noi
for any  $n \in \Z$, 
where $\ld = 
\max( 1- \al, 0)$ when $\al\ne 1$ and $\ld = \eps$ when $\al = 1$ for any $\eps > 0$.


\end{lemma}










\section{Modulated KdV}
\label{SEC:KDV}

We recall the following boundedness 
and nonlinear smoothing property 
of the bilinear operator $X^\KDV_{t, r}$;
see \cite[Proposition 4.1]{CGLLO1}.




\begin{lemma}\label{LEM:nonlin1}
Given $\rho \ge\frac 12$,  $\frac12< \g < 1$, and $T> 0$, 
let  $w$ be $(\rho,\g)$-irregular on $[0, T]$ in the sense of Definition~\ref{DEF:ir}.
Given $0 \le r < t \le T$, let $X^\KDV_{t, r}$
be as in \eqref{K1}.
Suppose that $s \in \R$ satisfies \eqref{nonlin1b}.
Then, 
we have 
$X^\KDV_{t, r} \in \L_2(H^s(\T); H^{s}(\T))$
with the following bound\textup{:}
\begin{align}
 \|X^\KDV_{t, r}\|_{\L_2(H^s; H^{s})}
\les 
(t- r)^\g  \|\Phi^w\|_{  \W^{\rho,\g}_T} 
\label{non1c}
\end{align}

\noi
for any $0 \le r < t \le T$, 
where the implicit constant is independent of $T > 0$.
In addition to~\eqref{nonlin1b}, suppose that $s_0 > s$ satisfies
one of the following conditions\textup{:}
\begin{align}
\begin{split}
\textup{(ii.a)} &\ \  
0 \le s + \rho \le  \tfrac 12
 \quad \text{and} \quad 
s_0 < 2s +  3\rho - \tfrac 32, \\
\textup{(ii.b)} &\ \  
s + \rho >  \tfrac 12 
\quad  \text{and} \quad s_0 \le s + 2\rho - 1.
\end{split}
\label{non1a}
\end{align}




\noi
Then, 
we have 
$X^\KDV_{t, r} \in \L_2(H^s(\T); H^{s_0}(\T))$
with the following bound\textup{:}
\begin{align*}
 \|X^\KDV_{t, r}\|_{\L_2(H^s; H^{s_0})}
\les 
(t-r)^\g  \|\Phi^w\|_{  \W^{\rho,\g}_T}
\end{align*}

\noi
for any $0 \le r < t \le T$.



\end{lemma}

From \eqref{non1a}, we see that, 
given any $s, s_0 \in \R$ with $s_0 \ge s$
(even if $s_0$ is much larger than~$s$), 
the bilinear operator $X^\KDV_{t, r}$
is bounded 
 from $H^s(\T) \times H^s(\T)$ into $H^{s_0}(\T)$, 
 provided that $\rho$ is sufficiently large.











We have the following boundedness property for
the trilinear operator  $\NN^\KDV_{t, r}$
defined in~\eqref{nox1}.

\begin{lemma}\label{LEM:nonlin2}
Given $\rho \ge\frac 12$,  $\frac12< \g < 1$, and $T> 0$, 
let  $w$ be $(\rho,\g)$-irregular on $[0, T]$ in the sense of Definition~\ref{DEF:ir}.
Given $0 \le r < t \le T$, let $\NN^\KDV_{t, r}$
be as in \eqref{nox1}.
Suppose that $s\ge 0$ and $s_0 \in \R$ satisfy
\begin{align}
s + \rho \ge  1
\label{nox1a}
\end{align}

\noi
and 
\begin{align}
s_0 < s + 2\rho - \frac 52.
\label{nox1b}
\end{align}

\noi
Then, 
$\NN^\KDV_{t, r}$  belongs to $\L_3(H^s(\T); H^{s_0}(\T))$, 
satisfying
 the following bound\textup{:}
\begin{align}
 \|\NN^\KDV_{t, r}\|_{\L_3(H^s; H^{s_0})}
\les 
(t- r)^\g  \|\Phi^w\|_{  \W^{\rho,\g}_T}
\label{nox1c}
\end{align}

\noi
for any $0 \le r < t \le T$, 
where the implicit constant is independent of $T > 0$.


\end{lemma}





%We now present a proof of  Lemma \ref{LEM:nonlin2}.


\begin{proof}%[Proof of Lemma \ref{LEM:nonlin2}]

From \eqref{nox1} with \eqref{rho1} and \eqref{K3a}, we have
\begin{align}
 \| \NN^\KDV_{t, r}(f_1, f_2, f_3)\|_{H^{s_0}}
&  \le
(t- r)^\g \|\Phi^w\|_{  \W^{\rho,\g}_T}M^\KDV(f_1, f_2, f_3)
\label{nox2}
\end{align}

\noi
for any $0 \le r < t \le T$, 
where $M^\KDV(f_1, f_2, f_3)$ is given by 
\begin{align}
\begin{split}
& M^\KDV(f_1, f_2, f_3)\\
& \ \  =\bigg\|
\sum_{\substack{n_1, n_2, n_3 \in \Z^*\\n = n_{123}}}
\frac{1}{\jb{n}^{\rho-s_0 - 1}
\jb{n_{12}}^{\rho-1}
\jb{n_1}^{s}\jb{n_2}^{s}\jb{n_3}^{\rho+s}} 
\prod_{j = 1}^3\jb{n_j}^s|\ft f_j(n_j)|
\bigg\|_{\l^2_n}.
\end{split}
\label{nox3}
\end{align}




From \eqref{nox3}, the triangle inequality $\jb{n_{12}}^s \les \jb{n_1}^s\jb{n_2}^s$
for $s \ge 0$, and Cauchy-Schwarz's inequality, we have 
\begin{align}
\begin{split}
 M^\KDV(f_1, f_2, f_3)
&   \le \bigg\|
\sum_{ n_3 \in \Z}
\frac{1}{\jb{n}^{\rho-s_0 - 1}
\jb{n - n_3}^{\rho-1+s}
\jb{n_3}^{\rho+s}} 
\jb{n_3}^s|\ft f_j(n_3)|\bigg\|_{\l^2_n}\\
 & \hphantom{XX} \times 
\sup_{n \in Z^*}\sum_{ n_1 \in \Z}
\ind_{n = n_{123}}
\prod_{j = 1}^2
\jb{n_j}^s|\ft f_j(n_j)|
\\
& 
\le 
 J^\KDV\cdot \prod_{j = 1}^3\|f_j\|_{H^s}, 
\end{split}
\label{nox4}
\end{align}


\noi
where  $ J^\KDV$ is given by 
\begin{align*}
 J^\KDV
= \bigg(\sum_{n,  n_3 \in \Z}
\frac{1}{\jb{n}^{2\rho-2s_0 - 2}
\jb{n - n_3}^{2\rho-2+2s}
\jb{n_3}^{2\rho+2s}} 
\bigg)^\frac 12.
\end{align*}

\noi
By
Lemma \ref{LEM:SUM}
in summing first in $n_3$ and then in $n$, we have $J^\KDV < \infty$, 
\noi
provided that 
$s + \rho  \ge 1$ and $s_0 < s + 2\rho - \frac 52$.
Then, the desired bound \eqref{nox1c} follows
from \eqref{nox2} and \eqref{nox4},  
provided that \eqref{nox1a}
and \eqref{nox1b}
hold.
\end{proof}



\begin{remark}\label{REM:non1}\rm
(i) By considering the contribution from 
the case $|n_1|\sim |n_2|\gg |n|, |n_3|, |n_{12}|$
in \eqref{nox1}, 
it is easy to see that the 
bound \eqref{nox1c} does not hold for $s < 0$.


\smallskip

\noi
(ii) 
By  the so-called Cauchy-Schwarz's argument
(see the proof of \cite[Proposition 4.1]{CGLLO1}), 
\begin{align*}
 \|\NN^\KDV_{t, r}(f_1, f_2, f_3)\|_{H^{s_0}}
 \le
(t-r)^\g  \|\Phi^w\|_{  \W^{\rho,\g}_T} 
\prod_{j = 1}^3\|f_j\|_{H^s}
\cdot \min_{j = 1, \dots 4}\sup_{n_j \in \Z^*} J_{n_j}, 
\end{align*}


\noi
where (by setting $n_4 = -n$) $J_{n_j}$ is given by 
\begin{align*}
J_{n_j} 
& = \bigg(
\sum_{\substack{n_1, n_2, n_3, n_4 \in \Z^*\\ n_{1234}= 0}}
\frac{\ind_{n_{12}\ne 0}}{\jb{n_4}^{2(\rho-s_0 - 1)}
\jb{n_{12}}^{2(\rho-1)}
\jb{n_1}^{2s}\jb{n_2}^{2s}\jb{n_3}^{2(\rho+s)}} 
 \bigg)^\frac{1}{2}.
\end{align*}


\noi
Then, we  bound $J_{n_j}$, using Lemma \ref{LEM:SUM}.
When $j = 3, 4$, 
this argument yields worse conditions than \eqref{nox1a} and \eqref{nox1b}.
When $j = 1, 2$, 
this argument also provides the conditions~\eqref{nox1a} and~\eqref{nox1b}.




\end{remark}












We conclude this section by  presenting a proof of Theorem \ref{THM:UU1}.

\begin{proof}[Proof of Theorem \ref{THM:UU1}]
(i) Let $\G(u) = \G_{u_0}(u)$ be the right-hand side of \eqref{NF4a}.
Under the hypothesis of Theorem~\ref{THM:UU1}\,(i), 
it follows from 
 Lemmas \ref{LEM:nonlin1} and 
\ref{LEM:nonlin2}
that 
\begin{align*}
& \| \G(u) \|_{C_\tau H^s_x}
 \le \| u_0\|_{H^s}
+ C\tau^\g
  \|\Phi^w\|_{  \W^{\rho,\g}_T} 
  \Big(\|u\|_{C_\tau H^s_x}^2 + \|u\|_{C_\tau H^s_x}^3\Big),\\
& \| \G(u)  - \G(v) \|_{C_\tau  H^s_x}\\
& \quad \le 
 C\tau^\g
  \|\Phi^w\|_{  \W^{\rho,\g}_T} 
\Big(1+ \|u\|_{C_\tau H^s_x}^2 + \|v\|_{C_\tau H^s_x}^2\Big) 
\|u- v \|_{C_\tau H^s_x}.
\end{align*}


\noi
for any $0 < \tau \le \min(T, 1)$.
Then, 
 a standard contraction argument
yields unconditional  local well-posedness
of \eqref{NF4a}
with a unique solution $u \in C([0, \tau]; H^s(\T))$, 
where the local existence time 
$\tau > 0$
depends only on 
$\|u_0\|_{H^s}$ and   $\|\Phi^w\|_{  \W^{\rho,\g}_T} $.
Here, the uniqueness of the solution a priori holds only in the ball 
$B_R \subset C([0, \tau]; H^s(\T))$
of radius $R \sim \|u_0\|_{H^s}$ centered at the origin
but, by a standard continuity argument, 
we can extend the uniqueness to the entire space $C([0, \tau]; H^s(\T))$
(by possibly making $\tau$ smaller by a multiplicative constant).
Moreover, by a standard persistence-of-regularity argument, 
we can easily show that the local existence time $\tau$
in fact depends only on 
$\|u_0\|_{L^2}$ and   $\|\Phi^w\|_{  \W^{\rho,\g}_T} $.
Since the argument is standard, we omit details.\footnote{See \cite[Proposition 4.1\,(i)]{CGLLO1}
for the persistence-of-regularity property for the bilinear operator $X_{t, 0}$ defined in \eqref{K1}.
For the trilinear operator $\NN_t$ defined in \eqref{nox1}, 
we can simply repeat the proof of Lemma \ref{LEM:nonlin2}
with $s_0 = s = 0$
by inserting $\jb{n}^s \les \jb{n_1}^s + \jb{n_2}^s+ \jb{n_3}^s$.
}




Next, we show that the solution $u$ to the normal form equation \eqref{NF4a}
constructed above 
is indeed a solution to the original modulated KdV \eqref{kdv1}.
From the contraction argument above, 
we also have the following Lipschitz bound:
\begin{align}
\| u^{(1)} - u^{(2)}\|_{C_\tau H^s_x}
\les \| u_0^{(1)} - u_0^{(2)}\|_{H^s}
\label{NF4c}
\end{align}

\noi
for any solution $u^{(j)}$  to the normal form equation \eqref{NF4a}
with initial data $u_0^{(j)}$, $j = 1, 2$, 
constructed by the contraction argument described above, 
where $\tau > 0$ denotes the minimum of the 
local existence time for $u^{(j)}$, $j = 1, 2$.
Now, given $u_0 \in H^s(\T)$ and
 a $(\rho,\g)$-irregular function $w$ on $[0, T]$, 
let $u \in C([0, \tau]; H^s(\T))$ be the solution  \eqref{NF4a}
with $u|_{t = 0} = u_0$, and 
$v \in C^\g([0, T]; H^s(\T))$ be the (global-in-time) solution to the 
modulated KdV \eqref{kdv1} with $v|_{t = 0} = u_0$ constructed in \cite{CGLLO1}.
By repeating the computations presented in Subsection~\ref{SUBSEC:NF}
(see also the discussion below for the justification
of the formal steps in~\eqref{NF2}
for functions in $C([0, T]; L^2(\T))$), 
we see that $v$ also satisfies the normal form equation~\eqref{NF4a}
with initial data $u_0$.
Therefore, from \eqref{NF4c}, we conclude that $v \equiv u$ on the time interval $[0, \tau]$
and thus that $u$ satisfies the original modulated KdV \eqref{kdv1}
(more precisely, the Duhamel formulation~\eqref{mild1x}).





It remains to  justify the formal steps in \eqref{NF2}.
For each fixed  $n \in \Z^*$, we need to justify
the following three steps:

\begin{itemize}

\smallskip

\item[(a)] 
switching of the time integration
and the summation at the second equality in \eqref{NF2}, 


\smallskip

\item[(b)] application of the product rule
on $\ft \uu(t', n_1)\ft \uu(t', n_2)$ 
at the third equality 
in \eqref{NF2}, 



\smallskip

\item[(c)] 
integration by parts
at the third equality 
in~\eqref{NF2}.

\end{itemize}

\smallskip


\noi
Fix  $u \in C([0, \tau]; L^2(\T))$ for some $0 < \tau \le T$.
(a)~Since we have
\begin{align*}
\sup_{0 \le t \le \tau} \sum_{\substack{n_1, n_2 \in \Z^*\\n = n_1 + n_2}}
|e^{i  \Xi_\KDV (\bar n)w(t')} \ft \uu(t', n_1)\ft \uu(t', n_2) |
\le \| \uu \|_{C_\tau L^2_x}^2
\end{align*}

\noi
for each fixed  $n \in \Z^*$,
the second equality of \eqref{NF2}
follows from 
the Fubini  theorem.
(b)~From~\eqref{kdv1b}, we have 
\begin{align*}
|\dt \ft \uu (t, n)| \le |n| \| u (t) \|_{L^2_x}^2
\end{align*}

\noi
for any $0 \le t \le T$.
In particular, 
for each fixed  $n \in \Z^*$,
we have $\ft \uu (\cdot, n)$
is a $C^1$-function in $t$, 
which justifies the application of 
the product rule 
at the third step in \eqref{NF2}.
(c)~Furthermore, by noting from~\eqref{rho2} that 
$\Phi^w_{t_1, t_2}(\Xi_\KDV (\bar n))$
is a $C^1$-function in $t_2$, 
the 
integration by parts
at the third equality 
in~\eqref{NF2}
is justified.
The discussion above shows that 
any solution $u \in C([0, \tau]; L^2(\T))$ to 
the Duhamel formulation \eqref{mild1x} of the modulated KdV \eqref{kdv1}
with $u|_{t = 0} = u_0$ for some $0 < \tau \le T$
(which does not have to be the solution constructed above
or that constructed in \cite{CGLLO1})
is also a solution the normal form equation~\eqref{NF4a}, 
which proves unconditional uniqueness
in the entire space $C([0, \tau]; L^2(\T))$
without intersecting with any auxiliary space.


Lastly, suppose in addition that 
 $w$ is  $(\rho,\g)$-irregular  on $\R_+$.
Then, the claimed unconditional global well-posedness of
the modulated KdV \eqref{kdv1}
in $H^s(\T)$ follows from 
the global well-posedness proven in \cite{CGLLO1}
and the local-in-time unconditional uniqueness
discussed above which can be iteratively applied
to yield global-in-time unconditional uniqueness.


\medskip

\noi
(ii) 
The claimed nonlinear smoothing follows
from applying
 Lemmas \ref{LEM:nonlin1} and 
\ref{LEM:nonlin2}
to the normal form equation \eqref{NF4a}. 
\end{proof}






\section{Modulated BO}
\label{SEC:BO}


We present a proof of Theorem \ref{THM:UU2}
for  the modulated BO \eqref{BO} on $\T$.
The bilinear driver associated with 
the modulated BO~\eqref{BO} on $\T$ is given by 
\begin{align}
X^\BO_{t,r}( f _1, f _2)
=\int_r^t  \uw  (t')^{-1}
\dx \big( (\uw (t')   f _1)( \uw (t')  f_2)\big) dt',
\label{BX1}
\end{align}


\noi
where 
\begin{align*}
\uw (t)=e^{w(t) \H\dx^2 }
\end{align*}

\noi
denotes 
the modulated linear propagator
for \eqref{BO}.
By taking the Fourier transform, we have 
\begin{align*}
\F\big( X^\BO_{t,r}  ( f _1, f _2) \big)(n)  =
in  \sum_{ \sub {n_1, n_2 \in \Z^*\\ n = n_1+n_2}}
 \Phi_{t,r}^{w}(\Xi_\BO (\bar n))\ft  f_1 ( n_1 )  \ft   f_2 (n_2), 
\end{align*}


\noi
where $\Phi_{t,r}^w$ is as in  \eqref{rho2}
and 
$\Xi_\BO (\bar n)$ denotes the resonance function for BO given by 
\begin{align}
\Xi_\BO (\bar n)
= \Xi_\BO(n,n_1,n_2) = - |n| n +  | n_1 |n_1 +  |n_2| n_2.
\label{BX4}
\end{align}



\noi
Given $0 \le r< t  \le T$, define the trilinear operator  $\NN^\BO_{t, r}(f_1, f_2, f_3)$, acting on functions on $\T$, 
 by 
\begin{align}
\begin{split}
& \F_x\big(\NN^\BO_{t, r}(f_1, f_2, f_3)\big)(n)\\
& \quad =   - 
2 n  \sum_{\substack{n_1, n_2, n_3 \in \Z^*\\n = n_{123}}}
 \Phi^w_{t, r}(\Xi_\BO (n, n_{12}, n_3))
%\\
%& \hphantom{XXX}
%\times 
e^{i  \Xi_\BO (n_{12}, n_1, n_2)w(r)}n_{12} 
\prod_{j = 1}^3\ft f_j(n_j), 
\end{split} 
\label{BX5}
\end{align}


\noi
where we used the convention \eqref{no1}.
Let $\uu(t) = \uw(t)^{-1} u(t)$ denote the 
modulated interaction representation 
of a solution $u$ to  (the Duhamel formulation of)~\eqref{BO}.
Then, 
proceeding as in Subsection \ref{SUBSEC:NF}, 
we have the following normal form equation for 
the 
modulated BO \eqref{BO}:
\begin{align*}
\ft \uu(t, n) - \ft \uu(0, n) 
 = \F_x\big(X^\BO_{t, 0}(\uu(0))\big)(n)
 +  \int_0^t 
\F_x\big(\NN^\BO_{t, t'}(\uu(t'))\big)(n)dt', 
\end{align*}

\noi
where we used the short-hand notation \eqref{short1}.


We recall the following boundedness 
and nonlinear smoothing property 
of the bilinear operator $X^\BO_{t, r}$;
see \cite[Proposition 5.1]{CGLLO1}.











\begin{lemma}\label{LEM:nonlin3}
Given $\rho \ge 1$,  $\frac12< \g < 1$, and $T> 0$, 
let  $w$ be $(\rho,\g)$-irregular on $[0, T]$ in the sense of Definition~\ref{DEF:ir}.
Given $0 \le r < t \le T$, let $X^\BO_{t, r}$
be as in \eqref{BX1}.
Suppose that $s \in \R$ satisfies~\eqref{regBO1}.
Then, 
we have $X_{t, r}^\BO \in \L_2(H^s(\T); H^{s}(\T))$
with the following bound\textup{:}
\begin{align*}
 \|X^\BO_{t, r}\|_{\L_2(H^s; H^{s})}
\les 
(t-r)^\g  \|\Phi^w\|_{  \W^{\rho,\g}_T}
\end{align*}

\noi
for any $0 \le r < t \le T$, 
where the implicit constant is independent of $T > 0$.
In addition to~\eqref{regBO1}, suppose that 
$\rho  > 1$ and $s_0 > s$ satisfy
one of the following conditions\textup{:}
\begin{align*}
\textup{(ii.a)} &\ \
0 \le s < s_0 \le \rho - 1,\\
\textup{(ii.b)} &\ \
 s <0
\quad 
\text{and}\quad 
s_0 \le s + \rho - 1.
\end{align*}




\noi
Then, 
we have 
$X^\BO_{t, 0} \in \L_2(H^s(\T); H^{s_0}(\T))$
with the following bound\textup{:}
\begin{align*}
 \|X^\BO_{t, r}\|_{\L_2(H^s; H^{s_0})}
\les 
(t- r)^\g  \|\Phi^w\|_{  \W^{\rho,\g}_T}
\end{align*}

\noi
for any $0 \le r < t \le T$.



\end{lemma}



Once we prove the following lemma
on the trilinear operator $\NN_{t, r}^\BO$
defined in \eqref{BX5},
Theorem~\ref{THM:UU2} 
for the modulated BO \eqref{BO} follows
from proceeding as in 
the proof of Theorem~\ref{THM:UU1}
for the modulated KdV \eqref{kdv1}
presented in Section \ref{SEC:KDV}.
We omit details of the proof of 
Theorem~\ref{THM:UU2}.


\begin{lemma}\label{LEM:nonlin4}
Given $\rho \ge1$,  $\frac12< \g < 1$, and $T> 0$, 
let  $w$ be $(\rho,\g)$-irregular on $[0, T]$ in the sense of Definition~\ref{DEF:ir}.
Given $0 \le r < t \le T$, let $\NN^\BO_{t, r}$
be as in \eqref{BX5}.
Suppose that $s\ge 0$ and $s_0 \in \R$ satisfy
\begin{align}
s + \rho \ge  1, \qquad 2s + \rho > \frac 32
\label{Bnox1a}
\end{align}

\noi
and 
\begin{align}
s_0 < s + \rho - \frac 52.
\label{Bnox1b}
\end{align}

\noi
Then, 
$\NN_{t, r}^\BO$  belongs to $\L_3(H^s(\T); H^{s_0}(\T))$, 
satisfying
 the following bound\textup{:}
\begin{align}
 \|\NN^\BO_{t, r}\|_{\L_3(H^s; H^{s_0})}
\les 
(t- r)^\g  \|\Phi^w\|_{  \W^{\rho,\g}_T}
\label{Bnox1c}
\end{align}

\noi
for any $0 \le r <  t \le T$, 
where the implicit constant is independent of $T > 0$.


\end{lemma}


Here, the condition $\rho \ge 1$ may be relaxed, 
but, in view of Lemma \ref{LEM:nonlin3}, 
 we keep it as it is since we need both Lemmas \ref{LEM:nonlin3}
 and \ref{LEM:nonlin4} to prove Theorem \ref{THM:UU2}.
A similar comment applies to Lemma \ref{LEM:nonlin6}.

Just as in the modulated KdV case, 
the bound \eqref{Bnox1c}
does not hold for $s < 0$; see
Remark~\ref{REM:non1}.
We conclude this section by  presenting a proof of  Lemma \ref{LEM:nonlin4}.


\begin{proof}[Proof of Lemma \ref{LEM:nonlin4}]





From \eqref{BX5} with \eqref{rho1} and \eqref{BX4}, we have
\begin{align*}
 \| \NN^\BO_{t, r}(f_1, f_2,  f_3)\|_{H^{s_0}}
&  \le
(t-r)^\g \|\Phi^w\|_{  \W^{\rho,\g}_T}M^\BO(f_1, f_2, f_3)
\end{align*}

\noi
for any $0 \le r < t \le T$, 
where $M^\BO(f_1, f_2, f_3)$ is given by 
\begin{align*}
& M^\BO(f_1, f_2, f_3)\\
& \ \  =\bigg\|
\sum_{\substack{n_1, n_2, n_3 \in \Z^*\\n = n_{123}}}
\frac{\jb{\Xi_\BO(n, n_{12}, n_3)}^{-\rho}}{\jb{n}^{-s_0 - 1}
\jb{n_{12}}^{-1}
\jb{n_1}^{s}\jb{n_2}^{s}\jb{n_3}^{s}} 
\prod_{j = 1}^3\jb{n_j}^s|\ft f_j(n_j)|
\bigg\|_{\l^2_n}.
\end{align*}


\noi
Then, arguing as in \eqref{nox4}, 
it suffices to show that 
\begin{align*}
 J^\BO
= \bigg(\sum_{n,  n_3 \in \Z^*}
\frac{\ind_{n = n_{123}}
\ind_{n_{12} \ne0}
\cdot \jb{\Xi_\BO(n, n_{12}, n_3)}^{-2\rho}}
{\jb{n}^{-2s_0 - 2}
\jb{n - n_3}^{-2+2s}
\jb{n_3}^{2s}} 
\bigg)^\frac 12 < \infty.
%\label{Bnox5}
\end{align*}







From \eqref{BX4}, we have, under $n = n_{12} + n_3$,  

\smallskip

\begin{itemize}
\item[(a)]
If $n_{12}n_3>0$ (i.e.~$n$, $n_{12}$, and $n_3$ all have the same sign), 
then  $|\Xi_\BO(n, n_{12}, n_3)|=2|n_{12}n_3|$.


\smallskip


\item[(b)]
If $n, n_{12}>0>n_3$,  then  $|\Xi_\BO(n, n_{12}, n_3)|=2|n n_3|$.



\smallskip


\item[(c)]


If $n, n_3>0>n_{12}$,  then  $|\Xi_\BO(n, n_{12}, n_3)|=2|n n_{12}|$.


\smallskip


\item[(d)]
If $n_{12}>0>n, n_3$,  then  $|\Xi_\BO(n, n_{12}, n_3)|=2|n n_{12}|$.



\smallskip


\item[(e)]
If $n_3>0>n, n_{12}$,  then  $|\Xi_\BO(n, n_{12}, n_3)|=2|n n_3|$.

\end{itemize}


\smallskip


\noi
By symmetry, 
we only consider Cases (a), (b),  and (c) in the following.



\medskip

\noi
$\bul$ {\bf Case (a):} $n_{12} n_3 >0$.\\
\indent
In this case, we have
\begin{align*}
( J^\BO)^2
\les  \sum_{n,  n_3 \in \Z}
\frac{1}
{\jb{n}^{-2s_0 - 2}
\jb{n - n_3}^{2\rho-2+2s}
\jb{n_3}^{2\rho+2s}} .
\end{align*}



\noi
By
Lemma \ref{LEM:SUM}
in summing first in $n_3$ and then in $n$, we have $J^\BO < \infty$, 
provided that 
$s + \rho  \ge 1$ and $s_0 < s + \rho - \frac 52$.












\medskip

\noi
$\bul$ {\bf Case (b):} $n, n_{12}>0>n_3$.\\
\indent
In this case, we have $|n_{12}| > |n_3|$ and thus $|n_{12}| \ges |n|$,
since $n = n_{12} + n_3$.
Then,  we have
\begin{align*}
 (J^\BO)^2
& \les  
\sum_{n,  n_3 \in \Z}
\frac{\ind_{n = n_{123}}\ind_{|n_{12}|\sim |n_3| \ges |n|}}
{\jb{n}^{2\rho-2s_0 - 2}
\jb{n_3}^{2\rho+4s - 2}} \\
& \quad + \sum_{n,  n_3 \in \Z}
\frac{\ind_{n = n_{123}}\ind_{|n_{12}|\sim |n| \ges |n_3|}}
{\jb{n}^{2\rho-2s_0 +2s- 4}
\jb{n_3}^{2\rho+2s}} 
< \infty, 
\end{align*}



\noi
provided that 
$2s + \rho  > \frac 32$ and $s_0 < s + \rho - \frac 52$.

\medskip

\noi
$\bul$ {\bf Case (c):} $n, n_3>0>n_{12}$.\\
\indent
In this case, we have $|n_{3}| > |n_{12}|$ and thus $|n_{3}| \ges |n|$,
since $n = n_{12} + n_3$.
Then, by Lemma~\ref{LEM:SUM},  we have
\begin{align*}
( J^\BO)^2
& \les  \sum_{n,  n_3 \in \Z}
\frac{\ind_{n = n_{123}}\ind_{|n_{12}|\sim |n_3| \ges |n|}}
{\jb{n}^{2\rho-2s_0 - 2}
\jb{n_3}^{2\rho + 4s-2}} \\
& 
\quad +   \sum_{n,  n_3 \in \Z}
\frac{\ind_{n = n_{123}}\ind_{|n|\sim |n_3| \ges |n_{12}|}}
{\jb{n}^{2\rho-2s_0 +2s - 2}
\jb{n - n_3}^{2\rho-2+2s}} < \infty, 
\end{align*}


\noi
provided that 
$2s + \rho  > \frac 32$ and $s_0 < \min\big(2s + 2\rho -3, s + \rho - \frac 32\big)$.




\medskip

Therefore, 
putting all the cases together, 
we conclude that  the desired bound \eqref{Bnox1c} holds, 
provided that \eqref{Bnox1a}
and \eqref{Bnox1b}
hold.
\end{proof}






\section{Modulated dNLS}
\label{SEC:NLS}


In this section, we briefly discuss a proof of Theorem \ref{THM:UU3}
for the modulated dNLS \eqref{dNLS1} on $\T$.
The bilinear driver associated with 
the modulated dNLS~\eqref{dNLS1} on $\T$ is given by 
\begin{align}
X^\dNLS_{t,r}( f _1, f _2)
=-i \int_r^t  \uw  (t')^{-1}
\dx \big( (\uw (t')   f _1)( \uw (t')  f_2)\big) dt',
\label{SX1}
\end{align}


\noi
where 
\begin{align*}
\uw (t)=e^{i w(t) \dx^2 }
\end{align*}

\noi
denotes 
the modulated linear propagator
for \eqref{dNLS1}.
By taking the Fourier transform, we have 
\begin{align*}
\F\big( X^\dNLS_{t,r}  ( f _1, f _2) \big)(n)  =
n  \sum_{ \sub {n_1, n_2 \in \Z^*\\ n = n_1+n_2}}
 \Phi_{t,r}^{w}(\Xi_\dNLS (\bar n))\ft  f_1 ( n_1 )  \ft   f_2 (n_2), 
\end{align*}


\noi
where $\Phi_{t,r}^w$ is as in  \eqref{rho2}
and 
$\Xi_\dNLS (\bar n)$ denotes the resonance function for dNLS given by 
\begin{align}
\begin{split}
\Xi_\dNLS (\bar n)
& = \Xi_\dNLS(n,n_1,n_2) = n^2 - n_1^2 - n_2^2\\
& =  2 n_1 n_2, 
\end{split}
\label{SX4a}
\end{align}

\noi
where the last equality holds under $n = n_1+ n_2$.
Given $0 \le r < t \le T$, 
define the trilinear operator $\NN^\dNLS_{t, r}(f_1, f_2, f_3)$, acting on functions on $\T$, 
 by 
\begin{align}
\begin{split}
& \F_x\big(\NN^\dNLS_{t, r}(f_1, f_2, f_3)\big)(n)\\
& \quad =   
- 2 i  n   \sum_{\substack{n_1, n_2, n_3 \in \Z^*\\n = n_{123}}}
 \Phi^w_{t, r}(\Xi_\dNLS (n, n_{12}, n_3))
%\\
%& \hphantom{XXX}
%\times 
e^{i  \Xi_\dNLS (n_{12}, n_1, n_2)w(r)}n_{12} 
\prod_{j = 1}^3\ft f_j(n_j), 
\end{split} 
\label{SX5}
\end{align}


\noi
where we used the convention \eqref{no1}.
Let $\uu(t) = \uw(t)^{-1} u(t)$ denote the 
modulated interaction representation 
of a solution $u$ to (the Duhamel formulation of)~\eqref{dNLS1}.
Then, 
proceeding as in Section \ref{SEC:KDV}, 
we have the following normal form equation for 
the 
modulated dNLS \eqref{dNLS1}:
\begin{align*}
\ft \uu(t, n) - \ft \uu(0, n) 
 = \F_x\big(X^\dNLS_{t, 0}(\uu(0))\big)(n)
 +  \int_0^t 
\F_x\big(\NN^\dNLS_{t, t'}(\uu(t'))\big)(n)dt', 
\end{align*}

\noi
where we used the short-hand notation \eqref{short1}.


The next lemma on the boundedness 
and nonlinear smoothing property 
of the bilinear operator $X^\dNLS_{t, r}$
follows from Cases (i.a) and (iii.a)
in the proof of 
\cite[Proposition 5.1]{CGLLO1}
for the bilinear driver $X^\BO_{t, r}$ 
for the modulated BO defined in~\eqref{BX1}.











\begin{lemma}\label{LEM:nonlin5}
Given $\rho \ge 1$,  $\frac12< \g < 1$, and $T> 0$, 
let  $w$ be $(\rho,\g)$-irregular on $[0, T]$ in the sense of Definition~\ref{DEF:ir}.
Given $0 \le r < t \le T$, let $X^\dNLS_{t, r}$
be as in \eqref{SX1}.
Suppose that $s \in \R$ satisfies 
\begin{align}
 s+ \rho \ge 0 \qquad\text{and}\qquad 
s >  \frac 32 - 2\rho.
%\begin{split}
%\textup{(i.a)} &\ \   1 \le \rho \le \frac 32  \quad \text{and} \quad s >  \frac 32 - 2\rho, \\
%\textup{(i.b)} &\ \  \rho >  \frac 32  \quad  \text{and} \quad  s \ge -  \rho.
%\end{split}
\label{regS1}
\end{align}



\noi
Then, 
we have $X_{t, r}^\dNLS \in \L_2(H^s(\T); H^{s}(\T))$
with the following bound\textup{:}
\begin{align*}
 \|X^\dNLS_{t, r}\|_{\L_2(H^s; H^{s})}
\les 
(t- r)^\g  \|\Phi^w\|_{  \W^{\rho,\g}_T}
\end{align*}

\noi
for any $0 \le r < t \le T$, 
where the implicit constant is independent of $T > 0$.
In addition to~\eqref{regS1}, suppose that 
$\rho  > 1$ and $s_0 > s$ satisfy
\begin{align*}
s+ \rho \ge 0, \qquad s_0 \le s + \rho - 1, \qquad \text{and}\qquad 
s_0 < 2s + 2\rho - \frac 32.
\end{align*}




\noi
Then, 
we have 
$X^\dNLS_{t, r} \in \L_2(H^s(\T); H^{s_0}(\T))$
with the following bound\textup{:}
\begin{align*}
 \|X^\dNLS_{t, r}\|_{\L_2(H^s; H^{s_0})}
\les 
(t- r)^\g  \|\Phi^w\|_{  \W^{\rho,\g}_T}
\end{align*}

\noi
for any $0 \le r< t \le T$.



\end{lemma}



Once we prove the following lemma
on the trilinear operator $\NN_{t, r}^\dNLS$
defined in \eqref{SX5}, 
Theorem~\ref{THM:UU3} 
for the modulated dNLS \eqref{dNLS1} follows
from proceeding as in 
the proof of Theorem~\ref{THM:UU1}
for the modulated KdV~\eqref{kdv1}
presented in Section \ref{SEC:KDV}.
We omit details of the proof of 
Theorem~\ref{THM:UU3}.


\begin{lemma}\label{LEM:nonlin6}
Given $\rho \ge 1$,  $\frac12< \g < 1$, and $T> 0$, 
let  $w$ be $(\rho,\g)$-irregular on $[0, T]$ in the sense of Definition~\ref{DEF:ir}.
Given $0 \le r < t \le T$, let $\NN^\dNLS_{t, r}$
be as in \eqref{SX5}.
Suppose that $s\ge 0$ and $s_0 \in \R$ satisfy
\begin{align}
s + \rho \ge  1 %, \qquad 2s + \rho > \frac 32
\label{Snox1a}
\end{align}

\noi
and 
\begin{align}
s_0 < s + \rho - \frac 52.
\label{Snox1b}
\end{align}

\noi
Then, 
$\NN_{t, r}^\dNLS$  belongs to $\L_3(H^s(\T); H^{s_0}(\T))$, 
satisfying
 the following bound\textup{:}
\begin{align}
 \|\NN^\dNLS_{t, r}\|_{\L_3(H^s; H^{s_0})}
\les 
(t- r)^\g  \|\Phi^w\|_{  \W^{\rho,\g}_T}
\label{Snox1c}
\end{align}

\noi
for any $0 \le r < t \le T$, 
where the implicit constant is independent of $T > 0$.


\end{lemma}



Note that 
the bound \eqref{Snox1c}
does not hold for $s < 0$; see
Remark~\ref{REM:non1}.
We conclude this section by  presenting a proof of  Lemma \ref{LEM:nonlin6}.


\begin{proof}[Proof of Lemma \ref{LEM:nonlin6}]





From \eqref{SX5} with \eqref{rho1} and \eqref{BX4}, we have
\begin{align*}
 \| \NN^\dNLS_{t, r}(f_1, f_2, f_3)\|_{H^{s_0}}
&  \le
(t-r)^\g \|\Phi^w\|_{  \W^{\rho,\g}_T}M^\dNLS(f_1, f_2, f_3)
\end{align*}

\noi
for any $0 \le r < t \le T$, 
where $M^\dNLS(f_1, f_2, f_3)$ is given by 
\begin{align*}
& M^\dNLS(f_1, f_2, f_3)\\
& \ \  =\bigg\|
\sum_{\substack{n_1, n_2, n_3 \in \Z^*\\n = n_{123}}}
\frac{\jb{\Xi_\dNLS(n, n_{12}, n_3)}^{-\rho}}{\jb{n}^{-s_0 - 1}
\jb{n_{12}}^{-1}
\jb{n_1}^{s}\jb{n_2}^{s}\jb{n_3}^{s}} 
\prod_{j = 1}^3\jb{n_j}^s|\ft f_j(n_j)|
\bigg\|_{\l^2_n}.
\end{align*}


\noi
Then, arguing as in \eqref{nox4}, 
it suffices to show that 
\begin{align}
 J^\dNLS
= \bigg(\sum_{n,  n_3 \in \Z^*}
\frac{\ind_{n = n_{123}}
\ind_{n_{12} \ne0}
\cdot \jb{\Xi_\dNLS(n, n_{12}, n_3)}^{-2\rho}}
{\jb{n}^{-2s_0 - 2}
\jb{n - n_3}^{-2+2s}
\jb{n_3}^{2s}} 
\bigg)^\frac 12 < \infty.
\label{Snox5}
\end{align}







From \eqref{SX4a}, we have
\begin{align}
\Xi_\dNLS(n, n_{12}, n_3)
=  2 n_{12} n_3.
\label{SX4b}
\end{align}

\noi
 under $n = n_{12} + n_3$.
Then, from \eqref{Snox5}
and \eqref{SX4b}, we have
\begin{align*}
( J^\dNLS)^2
\les  \sum_{n,  n_3 \in \Z}
\frac{1}
{\jb{n}^{-2s_0 - 2}
\jb{n - n_3}^{2\rho-2+2s}
\jb{n_3}^{2\rho+2s}} .
\end{align*}



\noi
By
Lemma \ref{LEM:SUM}
in summing first in $n_3$ and then in $n$, we have $J^\dNLS < \infty$, 
provided that \eqref{Snox1a}
and \eqref{Snox1b} hold.
This proves 
\eqref{Snox1c}.
\end{proof}










\appendix

\section{Unconditional uniqueness for the  modulated cubic NLS 
on the circle}
\label{SEC:A}

In this appendix, we study  the 
 following modulated cubic NLS on $\T$:
\begin{align}
i \dt u + \dx^2 u \cdot \dt w= |u|^2u.
\label{xNLS1}
\end{align}


\noi
In \cite{CG1},
Chouk and the first author showed that 
if   $w$ is $(\rho,\g)$-irregular on $\R_+$
for some $\rho > \frac 12 $ (and $\frac 12 < \g < 1$), 
then \eqref{xNLS1} is  globally well-posed
in $L^2(\T)$; see \cite[Theorem~1.8]{CG1}.
%see
%\cite[Theorem 1.4 and Example 1\,(2) with Remark 1.2\,(ii) and (iii)]{Robert2}.


Given a solution $u \in C([0, T]; L^2(\T))$ to \eqref{xNLS1}, we
define the following invertible gauge transform:
\begin{align}
v(t) = e^{-2i \int_\T |u(t)|^2 dx_{_\T}}u(t).
\label{gauge1}
\end{align}

\noi
In view of the $L^2$-conservation (see \cite[Subsection 3.3]{CG1}), 
a direct computation shows that the gauged function $v$ satisfies 
the following renormalized modulated cubic NLS:
\begin{align}
\begin{split}
i \dt v + \dx^2 v \cdot \dt w
& = \Big(|v|^2 - 2 \int_\T|v|^2 dx_{_\T} \Big)v \\
& = \NN(v) + \RR(v).
\end{split}
\label{xNLS2}
\end{align}


\noi
Here,  $\NN(v)$ and $\RR(v)$ denote the non-resonant 
 and  (doubly)  resonant parts of the renormalized nonlinearity,
respectively, 
defined by 
\begin{align}
\begin{split}
\F_x\big(\NN(v_1,v_2,v_3)\big)(n) 
& = \sum_{\substack{n = n_{123}^*\\ n\neq n_{1},n_{3}}}
\ft v_1 (n_1) \cj{\ft v_2 (n_2)} \ft v_3 (n_3), \\
\F_x\big(\RR(v_1,v_2,v_3)\big) (n) 
& = - 
\ft v_1 (n) \cj{\ft v_2 (n)} \ft v_3 (n), 
\end{split}
\label{non1}
\end{align}

\noi
where   we used 
 the short-hand notation \eqref{short1}
and 
the following short-hand notation:
\begin{align}
n_{1\cdots k}^* = n_1 - n_2 + n_3 -  \cdots + n_k
\label{no2}
\end{align}

\noi
for $k \in 2\N + 1$.
The renormalization via the gauge transform \eqref{gauge1}
essentially removes the resonant contribution from $n = n_1$ or $n_3$
in the original  cubic nonlinearity (at the expense
of introducing the doubly resonant term $\RR(v)$),
which has been used widely in the study of the cubic NLS
in the low regularity setting;
see, for example, \cite{BO96, Christ, GH, CO, GO, FOW, OW1, OW}.
In particular, 
 with the renormalized nonlinearity 
on the right-hand side of~\eqref{xNLS2}, 
the resonance (i.e.~$\Xi_{\NLS}(\bar n) = 0$, 
where $\Xi_{\NLS}(\bar n)$ is as in~\eqref{Xi9})
occurs only when $n = n_1 = n_2 = n_3$.






Our main goal in this appendix is to 
implement a Poincar\'e-Dulac normal form reduction for \eqref{xNLS2}
and prove the following theorem.

\begin{theorem}\label{THM:UU4}
Given $\rho > \frac 23$,  $\frac12< \g < 1$, and $T> 0$, 
let  $w$ be $(\rho,\g)$-irregular on $[0, T]$ in the sense of Definition~\ref{DEF:ir}.
Then, 
given any    $s\ge \frac 16$, 
the modulated cubic NLS equation~\eqref{xNLS1}
on  $\T$
is unconditionally locally well-posed in $H^s(\T)$.
In particular, 
if, in addition, $w$ is  $(\rho,\g)$-irregular  on $\R_+$, 
then the modulated cubic NLS equation \eqref{xNLS1}
on  $\T$
is unconditionally globally well-posed in $H^s(\T)$.

\end{theorem}


In view of  the cubic nonlinearity in \eqref{xNLS1}
and the (essentially scaling invariant) embedding $H^\frac 16 (\T)\subset L^3(\T)$, 
the regularity threshold $s \ge \frac 16$ 
in Theorem \ref{THM:UU4} 
is sharp
(within the $L^2$-based Sobolev scale). 
%to make
%sense of the cubic nonlinearity 
%as a spatial distribution for each fixed time
%(without intersecting with an auxiliary function space).



\medskip

The remaining part of this appendix is devoted to a proof of Theorem \ref{THM:UU4}.
In view of the equivalence of \eqref{xNLS1} and \eqref{xNLS2}
in $L^2(\T)$, 
Theorem \ref{THM:UU4} follows once we prove the claim for 
the renormalized modulated cubic NLS~\eqref{xNLS2}.
Hence, we focus on \eqref{xNLS2} in the following.

Let $\vv(t) = \uw(t)^{-1} v(t)$ denote the 
modulated interaction representation 
of a solution $v$ to~\eqref{xNLS2}.
Then, from \eqref{non1}, we have 
\begin{align}
\begin{split}
\dt \ft \vv(t, n)
& = 
-i  
\sum_{\substack{n = n_{123}^*\\ n\neq n_{1},n_{3}}}
e^{i  \Xi_\NLS (\bar n)w(t)} 
\prod_{j = 1}^3 {\ft \vv}^*(t, n_j)\\
& \quad  +   i     |\ft \vv(t, n)|^2\ft \vv(t, n)
\end{split}
\label{xNLS2a}
\end{align}


\noi
for $n \in \Z$, 
where $\Xi_\NLS (\bar n)$ denotes the resonance function for
\eqref{xNLS2} defined in \eqref{Xi9}
and 	
\begin{align}
\ft \vv^*(t, n_j)
= \begin{cases}
\ft \vv(t, n_j), & \text{when $j$ is odd,}\rule[-3mm]{0pt}{0pt} \\
\cj{\ft \vv(t, n_j)}, & \text{when $j$ is even.}
\end{cases}
\label{gauge2}
\end{align}

\noi
Then,  we have 
\begin{align}
\begin{split}
\ft \vv(t, n) - \ft \vv(0, n) 
& = 
- i \int_0^t \F_x\big(\uw  (t')^{-1}\NN(\uw  (t')\vv(t'))\big) dt'\\
& \quad  - i \int_0^t \F_x\big(\uw  (t')^{-1}\RR(\uw  (t')\vv(t'))\big) dt'\\
&
= 
-i  
\int_0^t
\sum_{\substack{n = n_{123}^*\\ n\neq n_{1},n_{3}}}
e^{i  \Xi_\NLS (\bar n)w(t')} 
\prod_{j = 1}^3 \ft \vv^*(t', n_j)\\
& \quad  +   i   \int_0^t  |\ft \vv(t', n)|^2\ft \vv(t', n) dt'.
\end{split}
\label{xNLS3}
\end{align}


\noi
The first and second terms on the right-hand side of \eqref{xNLS3} represent 
the non-resonant and resonant  contributions, respectively.
In the following, we apply a normal form reduction
to the non-resonant part
(corresponding to 
a Poincar\'e-Dulac normal form reduction;
see Remark~\ref{REM:NF1}).




The trilinear  driver associated with
the non-resonant part in \eqref{xNLS3}  is given by 
\begin{align}
X^\NLS_{t,r}( f _1, f _2, f_3)
=-i \int_r^t  \uw  (t')^{-1}
\NN \big( \uw (t')   f _1, \uw (t')  f_2, 
 \uw (t')  f_3
 \big) dt',
\label{xX1}
\end{align}


\noi
where 
\begin{align*}
\uw (t)=e^{i w(t) \dx^2 }
\end{align*}

\noi
denotes 
the modulated linear propagator
for \eqref{xNLS2}.
By taking the Fourier transform, 
it follows from \eqref{non1} that  
\begin{align}
\F\big( X^\NLS_{t,r}  ( f _1, f _2, f_3) \big)(n)  =
-i  \sum_{ \sub {n = n_{123}^*\\n \ne n_1, n_3}}
 \Phi_{t,r}^{w}(\Xi_\NLS (\bar n))\ft  f_1 ( n_1 ) \cj{ \ft   f_2 (n_2)}
 \ft  f_3 ( n_3 ), 
\label{xX2}
\end{align}


\noi
where $\Phi_{t,r}^w$ 
and 
$\Xi_\NLS (\bar n)$ 
are as in  \eqref{rho2}
and  \eqref{Xi9}.





Proceeding as in \eqref{NF2} with \eqref{xNLS3} and \eqref{NS1}, we have 
\begin{align}
\begin{split}
- i  & \int_0^t \F_x\big(\uw(t')^{-1}\NN(\uw(t') \vv(t')\big) dt'\\
& =  i  
\sum_{\substack{n = n_{123}^*\\ n\neq n_{1},n_{3}}}
 \int_0^t
 \partial_{t'} \Phi^w_{t, t'}(\Xi_\NLS (\bar n))
 \prod_{j = 1}^3 \ft \vv^*(t', n_j)
 dt'\\
& = \F_x\big(X^\NLS_{t, 0}(\vv(0), \vv(0), \vv(0))\big)(n)  \\
%\\
&  \quad -  i \sum_{\substack{n = n_{123}^*\\ n\neq n_{1},n_{3}}}
 \int_0^t
 \Phi^w_{t, t'}(\Xi_\NLS (\bar n))
\, \dt \bigg( \prod_{j = 1}^3 \ft \vv^*(t', n_j)\bigg) 
 dt'.
 \end{split}
\label{xNLS4}
\end{align}


\noi
In view of the second term on the right-hand side of \eqref{xNLS4}
with \eqref{xNLS2a}, 
we define the 
quintilinear operators  
$\NN \NN^\NLS_{t, r}(f_1, \dots, f_5)$
and 
$\NN \RR^\NLS_{t, r}(f_1, \dots, f_5)$, acting on functions on $\T$, 
 by 
\begin{align}
\begin{split}
 \F_x & \big(\NN \NN^\NLS_{t, r}(f_1, \dots, f_5)\big)(n)\\
&  =   - 
2   \sum_{\substack{n = n_{12345}^*\\
n \ne n_{123}^*, n_5\\
n_{123}^* \ne n_1, n_3
}}
 \Phi^w_{t, r}(\Xi_\NLS (n, n_{123}^*, n_4, n_5))
e^{i  \Xi_\NLS (n_{123}^*, n_1, n_2, n_3)w(r)}
\prod_{j = 1}^5{\ft f_j}^*(n_j)\\
& \quad  + 
   \sum_{\substack{n = n_{12345}^*\\
n \ne n_1, n_5\\
n_{234}^* \ne n_2, n_4
}}
 \Phi^w_{t, r}(\Xi_\NLS (n, n_1, n_{234}^*,  n_5))
e^{- i  \Xi_\NLS (n_{234}^*, n_2, n_3, n_4)w(r)}
\prod_{j = 1}^5{\ft f_j}^*(n_j)\\
& =: \F_x\big(\1_{t, r}(f_1, \dots, f_5)\big)(n) + \F_x\big(\II_{t, r}(f_1, \dots, f_5)\big)(n)
\end{split} 
\label{xX5}
\end{align}

\noi
and
\begin{align}
\begin{split}
 \F_x & \big(\NN \RR^\NLS_{t, r}(f_1, \dots, f_5)\big)(n)\\
&  =    
2   
\sum_{\substack{n = n_1 - n_4 + n_5\\ n\neq n_{1},n_{5}}}
 \Phi^w_{t, r}(\Xi_\NLS (n, n_1, n_4, n_5))
\prod_{j = 1}^3{\ft f_j}^*(n_1)
\cdot \prod_{j = 4}^5{\ft f_j}^*(n_j)\\
& \quad -
\sum_{\substack{n = n_1 - n_2 + n_5\\ n\neq n_{1},n_{5}}}
 \Phi^w_{t, r}(\Xi_\NLS (n, n_1, n_2, n_5))
\prod_{j = 2}^4{\ft f_j}^*(n_2)
\cdot \prod_{j \in \{1, 5\}} {\ft f_j}^*(n_j)
\end{split} 
\label{xX6}
\end{align}



\noi
for $0 \le r< t  \le T$, 
where we used the conventions \eqref{no2}
and \eqref{gauge2}.

Hence, putting 
\eqref{xNLS3}, \eqref{xNLS4}, \eqref{xX5}, and \eqref{xX6}
together, we obtain
 the following normal form equation for 
the 
renormalized modulated cubic NLS \eqref{xNLS2}:
\begin{align}
\begin{split}
\ft \vv(t, n) - \ft \vv(0, n) 
& =    \F_x\big(X^\NLS_{t, 0}(\vv(0))\big)(n)  
 +    \int_0^t \F_x\big(\Rf  (\vv(t'))\big) dt'\\
& \quad  +  \int_0^t 
 \F_x\big(\NN\NN^\NLS_{t, t'}(\vv(t'))\big)(n)dt'\\
& \quad   +  \int_0^t 
 \F_x\big(\NN\RR^\NLS_{t, t'}(\vv(t'))\big)(n)dt', 
\end{split}
\label{xNLS5}
\end{align}


\noi
where we used the short-hand notation \eqref{short1}
and the trilinear operator $\Rf(f_1, f_2, f_3)$, acting on functions on $\T$, 
is defined by 
\begin{align}
\F_x\big(\Rf(f_1, f_2, f_3)\big)(n)
= i \ft f_1 (n)\cj{\ft f_2 (n)}\ft f_3 (n).
\label{xNLS6}
\end{align}





We now state lemmas on boundedness properties
of the multilinear operators appearing in~\eqref{xNLS5}.
Once we prove these lemmas
(Lemmas \ref{LEM:nonlin7}, 
\ref{LEM:nonlin8}, and \ref{LEM:nonlin9}), 
Theorem~\ref{THM:UU4} 
for the renormalized modulated cubic NLS \eqref{xNLS2} 
(and hence for the modulated cubic NLS \eqref{xNLS1})
follows
from proceeding as in 
the proof of Theorem~\ref{THM:UU1}
for the modulated KdV~\eqref{kdv1}
presented in Section \ref{SEC:KDV}
with the observation that,   
given a solution  $\vv$
to~\eqref{xNLS2}
in the class
\[   C([0, \tau]; H^\frac 16(\T))
\subset C([0, \tau]; L^3(\T))
\]

\noi
for some $0 < \tau \le T$, 
 $\ft \vv (\cdot, n)$
is a $C^1$-function in $t$
for each fixed  $n \in \Z^*$, 
which is a key point in justifying the formal computations
in deriving the normal form equation \eqref{xNLS5};
see also~\cite{GKO, OW}.
We omit details of the proof of 
Theorem~\ref{THM:UU4}.


\medskip

The boundedness of $\Rf$ in \eqref{xNLS6} follows
trivially.


\begin{lemma}\label{LEM:nonlin7}
Let $\Rf$ be as in \eqref{xNLS6}.
Then, given  any 
 $ s \ge 0$, 
we have 
\begin{align}
 \|\Rf\|_{\L_3(H^s; H^{s})}
\le 1.
\label{non7}
\end{align}



\end{lemma}

\begin{proof}
The bound \eqref{non7} follows from $\jb{n}^s \le \jb{n}^{3s }$
for $s \ge 0$
and $\| a_n \|_{\l^6_n}\le \| a_n \|_{\l^2_n}$.
\end{proof}







Next, we establish boundedness of 
the trilinear operator  $X^\NLS_{t, r}$
defined in  \eqref{xX1}.











\begin{lemma}\label{LEM:nonlin8}
Given $\rho > \frac 12 $,  $\frac12< \g < 1$, and $T> 0$, 
let  $w$ be $(\rho,\g)$-irregular on $[0, T]$ in the sense of Definition~\ref{DEF:ir}.
Given $0 \le r < t \le T$, let $X^\NLS_{t, r}$
be as in \eqref{xX1}.
Then, 
given any $s \ge 0$, 
we have $X_{t, r}^\NLS \in \L_3(H^s(\T); H^{s}(\T))$
with the following bound\textup{:}
\begin{align*}
 \|X^\NLS_{t, r}\|_{\L_3(H^s; H^{s})}
\les 
(t-r)^\g  \|\Phi^w\|_{  \W^{\rho,\g}_T}
\end{align*}

\noi
for any $0 \le r < t \le T$, 
where the implicit constant is independent of $T > 0$.





\end{lemma}

\begin{proof}

From 
\eqref{xX2},  Cauchy-Schwarz's inequality, 
\eqref{rho1}, and \eqref{Xi9} (note that $n \ne n_1, n_3$)
with the triangle inequality
 $\jb{n_{123}^*}^s \les \jb{n_1}^s\jb{n_2}^s\jb{n_3}^s$
for $s \ge 0$, 
we have 
\begin{align*}
\| X^\NLS_{t,r}  ( f _1, f _2, f_3) \|_{H^s}
& \les 
(t - r)^\g
\|\Phi^w\|_{  \W^{\rho,\g}_T} 
\Bigg\|\bigg( \sum_{n_1, n_3 \in \Z}
\frac {1}{\jb{n - n_1}^{2\rho}\jb{n - n_3}^{2\rho}}
\bigg)^\frac 12\\
& \quad \times \bigg(  
\sum_{n = n_{123}^*}
\prod_{j = 1}^3
\jb{n_j}^{2s}
 |\ft  f_j ( n_j )|^2\bigg)^\frac 12 
 \Bigg\|_{\l^2_n}\\
 & \les
 \|\Phi^w\|_{  \W^{\rho,\g}_T} 
(t - r)^\g
\prod_{j = 1}^3 \|f_j\|_{H^s}, 
\end{align*}

\noi
provided that $\rho > \frac 12$.
\end{proof}




Lastly,  we establish boundedness of 
the quintilinear operators  
$\NN \NN^\NLS_{t, r}(f_1, \dots, f_5)$
and 
$\NN \RR^\NLS_{t, r}(f_1, \dots, f_5)$
defined in  \eqref{xX5}
and \eqref{xX6}.




\begin{lemma}\label{LEM:nonlin9}
\textup{(i)}
Given $\rho > \frac 12 $,  $\frac12< \g < 1$, and $T> 0$, 
let  $w$ be $(\rho,\g)$-irregular on $[0, T]$ in the sense of Definition~\ref{DEF:ir}.
Given $0 \le r < t \le T$, let $\NN \RR^\NLS_{t, r}$
be as in \eqref{xX6}.
Then, given  $s \ge 0$, 
$\NN \RR^\NLS_{t, r}$  belongs to $\L_5(H^s(\T); H^{s}(\T))$, 
satisfying
 the following bound\textup{:}
\begin{align}
 \|\NN \RR^\NLS_{t, r}\|_{\L_5(H^s; H^{s})}
\les 
(t- r)^\g  \|\Phi^w\|_{  \W^{\rho,\g}_T}.
\label{Xnox1c}
\end{align}




\medskip

\noi
\textup{(ii)}
Given $\rho > \frac 23 $,  $\frac12< \g < 1$, and $T> 0$, 
let  $w$ be $(\rho,\g)$-irregular on $[0, T]$ in the sense of Definition~\ref{DEF:ir}.
Given $0 \le r < t \le T$, let $\NN \NN^\NLS_{t, r}$
be as in \eqref{xX5}.
Then, given  $s \ge \frac 16$, 
$\NN \NN^\NLS_{t, r}$  belongs to $\L_5(H^s(\T); H^{s}(\T))$, 
satisfying
 the following bound\textup{:}
\begin{align}
 \|\NN \NN^\NLS_{t, r}\|_{\L_5(H^s; H^{s})}
\les 
(t- r)^\g  \|\Phi^w\|_{  \W^{\rho,\g}_T}.
\label{Xnox1d}
\end{align}



\end{lemma}







\begin{proof}
(i) 
From 
\eqref{xX6}, \eqref{xX2}, \eqref{xNLS6}, we have 
\begin{align*}
\NN \RR^\NLS_{t, r}(f_1, \dots, f_5)
= 2X^\NLS_{t,r}  ( \Rf(f_1, f_2, f_3) , f _4, f_5)
+ X^\NLS_{t,r}  ( f_1, \Rf(f_2, f_3, f_4) ,  f_5).
\end{align*}

\noi
Hence, the bound
\eqref{Xnox1c} follows from Lemmas \ref{LEM:nonlin7}
and \ref{LEM:nonlin8}, 
provided that $\rho > \frac 12 $ and $s \ge 0$.


\medskip





\noi
(ii) 
In view of the triangle inequality $\jb{n_{1\cdots5}^*}^\s \les \prod_{j= 1}^5 \jb{n_j}^\s$
for $\s \ge 0$, we only consider the case $s = \frac 16 $.


We first estimate $\1_{t, r}(f_1, \dots, f_5)$ in \eqref{xX5}.
\medskip


\noi
$\bul$ {\bf Case (a):} $|n|\gg |n_5|$.\\ %$|n_{123}^*| \gg |n_4|, |n_5|$.\\
\indent
From \eqref{xX5}, \eqref{rho1}, and \eqref{Xi9}
with 
$\jb{n} \sim \jb{n - n_5}$, 
we have 
\begin{align}
\begin{split}
 \|\1_{t, r}(f_1, \dots, f_5)\|_{H^\frac 16}
&  
\les 
(t - r)^\g
\|\Phi^w\|_{  \W^{\rho,\g}_T} 
\bigg\| \sum_{m = n_{123}^*}
\prod_{j = 1}^3
{\ft f_j}^*(n_j)\bigg\|_{\l^\infty_{m}}\\
& \quad \times \bigg\| 
\sum_{n = m - n_4 + n_5}
\frac {|\ft f_4(n_4)|  |\ft f_5(n_5)|}{\jb{n - m}^{\rho}\jb{n - n_5}^{\rho - \frac 16}}
\bigg\|_{\l^2_n}.
\end{split}
\label{non91}
\end{align}


\noi
By Hausdorff-Young's and Sobolev's inequalities, we have 
\begin{align}
\bigg\| \sum_{m = n_{123}^*}
\prod_{j = 1}^3
{\ft f_j}^*(n_j)\bigg\|_{\l^\infty_{m}}
\le \prod_{j = 1}^3 \|f_j\|_{L^3}\les \prod_{j = 1}^3 \|f_j\|_{H^\frac 16}.
\label{non92}
\end{align}


\noi
On the other hand, 
using  $\jb{n} \sim \jb{n - n_5}$
and applying Young's,  
Hausdorff-Young's,  and Sobolev's inequalities, we have 
the last factor on the right-hand side of \eqref{non91} is bounded by 
\begin{align}
\begin{split}
& \les 
\bigg\|\frac 1{\jb{n}^{\rho - \frac 16}}
\bigg\|_{\l^2_n}
\cdot \sum_{n_4, n_5 \in \Z}
\frac {|\ft f_4(n_4)|  |\ft f_5(n_5)|}{\jb{n_4 - n_5}^{\rho}}\\
& \les 
\| \ft f_4(n_4)\|_{\l^\frac 32_{n_4}}  
\bigg\|
\sum_{n_5 \in \Z}
\frac{\ft f_5(n_5)}{\jb{n_4 - n_5}^{\rho}}\bigg\|_{\l^3_{n_4}}\\
& \les \prod_{j = 4}^5 \|f_j\|_{L^3}\les \prod_{j = 4}^5 \|f_j\|_{H^\frac 16}, 
\end{split}
\label{non93}
\end{align}

\noi
provided that $\rho > \frac 23$.
Hence, the bound \eqref{Xnox1d}
follows from \eqref{non91}, 
\eqref{non92}, and 
\eqref{non93} in this case.


\medskip


\noi
$\bul$ {\bf Case (b):} $|n|\les |n_5|$.\\ %$|n_{123}^*| \gg |n_4|, |n_5|$.\\
\indent
By $\jb{n}^\frac 16 \les\jb{n_5}^\frac 16$, we have 
\begin{align}
\begin{split}
 \|\1_{t, r}(f_1, \dots, f_5)\|_{H^\frac 16}
&  
\les 
(t - r)^\g
\|\Phi^w\|_{  \W^{\rho,\g}_T} 
\bigg\| \sum_{m = n_{123}^*}
\prod_{j = 1}^3
{\ft f_j}^*(n_j)\bigg\|_{\l^\infty_{m}}\\
& \quad \times \bigg\| 
\sum_{n = m - n_4 + n_5}
\frac {|\ft f_4(n_4)|  \jb{n_5}^\frac 16 |\ft f_5(n_5)|}{\jb{n - m}^{\rho}\jb{n - n_5}^{\rho}}
\bigg\|_{\l^2_n}, 
\end{split}
\label{non94}
\end{align}

\noi
where the last factor is bounded by 
\begin{align}
\begin{split}
& \les 
\bigg\|
\sum_{n_4 \in \Z}
\frac{\ft f_4(n_4)}{\jb{n_4 - n_5}^{\rho}}\bigg\|_{\l^3_{n_5}}
\Bigg\|
\bigg\|
\frac{  \jb{n_5}^\frac 16 |\ft f_5(n_5)|}{\jb{n - n_5}^{\rho}}\bigg\|_{\l^\frac 32_{n_5}}
\Bigg\|_{\l^2_n}\\
& \les  \|f_4\|_{L^3}
\|f_5\|_{H^\frac 16}
\les \prod_{j = 4}^5 \|f_j\|_{H^\frac 16}, 
\end{split}
\label{non95}
\end{align}

\noi
provided that $\rho > \frac 23$.
Hence, the bound \eqref{Xnox1d}
follows from \eqref{non94}, 
\eqref{non92}, and 
\eqref{non95} in this case.


\medskip

Next, we estimate
the second term
 $\II_{t, r}(f_1, \dots, f_5)$ in \eqref{xX5}.



\medskip


\noi
$\bul$ {\bf Case (c):} $|n|\gg \max(|n_1|, |n_5|)$.\\ %$|n_{123}^*| \gg |n_4|, |n_5|$.\\
\indent
Proceeding as in \eqref{non91} with $\jb{n- n_5} \sim \jb{n}$, we have
\begin{align}
\begin{split}
 \|\II_{t, r}(f_1, \dots, f_5)\|_{H^\frac 16}
&  
\les 
(t - r)^\g
\|\Phi^w\|_{  \W^{\rho,\g}_T} 
\bigg\| \sum_{m = n_{234}^*}
\prod_{j = 2}^4
{\ft f_j}^*(n_j)\bigg\|_{\l^\infty_{m}}\\
& \quad \times \bigg\| 
\sum_{n = n_1 -m + n_5}
\frac {|\ft f_1(n_1)|  |\ft f_5(n_5)|}{\jb{n - n_1}^{\rho}\jb{n}^{\rho - \frac 16}}
\bigg\|_{\l^2_n}.
\end{split}
\label{non96}
\end{align}

\noi
With $\jb{n-n_1}\ges \jb{n_1 - n_5}$, we can bound the last factor 
of the right-hand side of \eqref{non96}
by~\eqref{non93} (after relabelling), 
while the penultimate factor of the right-hand side of \eqref{non96}
can be bounded by~\eqref{non92} (after relabelling).
This proves \eqref{Xnox1d} in this case, provided that $\rho > \frac 23$.


\medskip


\noi
$\bul$ {\bf Case (d):} $|n|\les \max(|n_1|, |n_5|)$.\\ %$|n_{123}^*| \gg |n_4|, |n_5|$.\\
\indent
Without loss of generality, assume that $|n_5|\ges |n_1|$.
Then, 
proceeding as in \eqref{non94},  we have 
\begin{align}
\begin{split}
 \|\II_{t, r}(f_1, \dots, f_5)\|_{H^\frac 16}
&  
\les 
(t - r)^\g
\|\Phi^w\|_{  \W^{\rho,\g}_T} 
\bigg\| \sum_{m = n_{234}^*}
\prod_{j = 2}^4
{\ft f_j}^*(n_j)\bigg\|_{\l^\infty_{m}}\\
& \quad \times \bigg\| 
\sum_{n = n_1-m  + n_5}
\frac {|\ft f_1(n_1)|  \jb{n_5}^\frac 16 |\ft f_5(n_5)|}{\jb{n - n_1}^{\rho}\jb{n - n_5}^{\rho}}
\bigg\|_{\l^2_n}, 
\end{split}
\label{non97}
\end{align}

\noi
where the last factor is bounded by 
\begin{align}
\begin{split}
& \les 
\bigg\|
\sum_{n_1 \in \Z}
\frac{\ft f_1(n_1)}{\jb{n - n_1}^{\rho}}\bigg\|_{\l^3_{n}}
\bigg\|
\sum_{n_5 \in \Z} \frac{  \jb{n_5}^\frac 16 |\ft f_5(n_5)|}{\jb{n - n_5}^{\rho}}\bigg\|_{\l^6_{n}}\\
& \les  \|f_1\|_{L^3}
\|f_5\|_{H^\frac 16}
\les \prod_{j \in \{1, 5\}} \|f_j\|_{H^\frac 16}, 
\end{split}
\label{non98}
\end{align}

\noi
provided that $\rho > \frac 23$.
Hence, the bound \eqref{Xnox1d}
follows from \eqref{non97}, 
and \eqref{non98}
with 
\eqref{non92} (after relabelling).
\end{proof}


















\begin{remark} \label{REM:NLS1}\rm



By slightly modifying the normal form reduction, 
 it is possible to prove nonlinear smoothing 
for \eqref{xNLS1},
as in Theorems \ref{THM:UU1}, \ref{THM:UU2}, and \ref{THM:UU3}.
More precisely, 
as  in 
\cite{GKO, OW}, 
we can implement a normal form reduction
%we can improve the regularity threshold
% at the level of the normal form equation
by introducing a suitable cutoff $K$ for the resonance function $\Xi_{\NLS}(\bar n)$ in \eqref{Xi9}
and thus separately considering
the {\it nearly resonant} contribution 
(coming from $|\Xi_{\NLS}(\bar n)|\ll K$)
and the {\it highly
non-resonant} contribution
(coming from $|\Xi_{\NLS}(\bar n)|\ges K$).
For example, take $K \sim n_{\max}$.
In  the nearly resonant case: $|\Xi_\NLS(\bar n)| \ll n_{\max}$, 
where 
$\Xi_\NLS(\bar n)$ and 
$n_{\max}$ are as in 
\eqref{Xi9} and \eqref{max1}, respectively, 
we have
\begin{align*}
 |n|\sim |n_1|\sim|n_2|\sim|n_3|,
\end{align*}

\noi
which implies nonlinear smoothing, 
provided that $s_0 \le 3s$
(possibly with further conditions).
On the other hand, 
it follows from~\eqref{rho1} 
and $|\Xi_\NLS(\bar n)| \ges  n_{\max}$
that  nonlinear smoothing of arbitrary degree
holds 
for the highly non-resonant part
provided that $\rho$ is sufficiently large.
For conciseness of the presentation, we omit details.


\end{remark}






\begin{remark}\rm

The regularity restriction $s \ge \frac 16$
appears in 
\eqref{non92}, 
which shows that if we were to stop normal form reductions
within a finite number of steps, 
we would not be able to go below 
the regularity threshold $ s= \frac 16$.
Hence, as in  the case of the unmodulated cubic NLS  \eqref{NLS3}
studied in 
 \cite{GKO, OW}, 
we expect that an infinite iteration of normal form reduction is needed
to go below $ s= \frac 16$.
We, however,  point out
that this comes with a twist due to 
the fact that the identity \eqref{NS0}
is replaced by 
\eqref{NS1} in the modulated setting.
We will address this issue in a forthcoming paper.




\end{remark}



\begin{ackno}\rm
The authors would like to thank Andreia Chapouto
for a comment on numerical schemes
in Remark \ref{REM:num}.
T.O.~was supported by the European Research Council (grant no.~864138 ``SingStochDispDyn")
and  by the EPSRC 
Mathematical Sciences
Small Grant  (grant no.~EP/Y033507/1).


\end{ackno}







\begin{thebibliography}{99}




\bibitem{ABLL}
S.~Athreya, O.~Butkovsky, K.~L\^e, L.~Mytnik,
{\it Well-posedness of stochastic heat equation with distributional drift and skew stochastic heat equation},
Comm. Pure Appl. Math. 77 (2024), no. 5, 2708--2777.


\bibitem{BIT}
A.V.~Babin, A.A.~Ilyin, E.S.~Titi,
{\it On the regularization mechanism for the periodic Korteweg-de Vries equation}, 
Comm. Pure Appl. Math. 64 (2011), no. 5, 591--648.



%\bibitem{BG24}
%D.~Bambusi, P.~G\'erard,
%{\it A Nekhoroshev theorem for some perturbations of the Benjamin-Ono equation with initial data close to finite gap tori}, 
%Math. Z. 307 (2024), no.3, Paper No. 54, 20 pp.
%


%
%
%\bibitem{Be66}
%T.B.~Benjamin,
%{\it  Internal waves of finite amplitude and permanent form}, 
%J. Fluid Mech. 25 (1966), 241--270. 
%




%
%
%\bibitem{Be67}
%T.B.~Benjamin,
%{\it  Internal waves of permanent form in fluids of great depth}, 
%J. Fluid Mech. 29 (1967), 559--592. 
%



\bibitem{BOP2}
\'A.~B\'enyi, T.~Oh, O.~Pocovnicu,
{\it On the probabilistic Cauchy theory of the cubic nonlinear Schr\"odinger equation on $\R^d$, $d\ge 3$}, 
Trans. Amer. Math. Soc. Ser. B 2 (2015), 1--50. 




%
%
%\bibitem{bessaih_evolution_2005}
%H.~Bessaih, M.~Gubinelli, F.~Russo,
%{\it The evolution of a random vortex filament}, 
%Ann. Probab. 33 (2005), no. 5, 1825--1855.
%





\bibitem{BO93}
J.~Bourgain, 
{\it Fourier transform restriction phenomena for certain lattice subsets and applications to nonlinear evolution equations, II. The KdV-equation,} 
 Geom. Funct. Anal. 3 (1993), no. 3, 209--262.


\bibitem{BO96}
J.~Bourgain,
{\it Invariant measures for the 2D-defocusing nonlinear Schr\"odinger equation},
Comm. Math. Phys. 176 (1996), no. 2, 421--445.

%
\bibitem{BO97}
J.~Bourgain, 
{\it Periodic Korteweg de Vries equation with measures as initial data}, 
 Selecta Math. (N.S.) 3 (1997), no. 2, 115--159.
%

%
%
%\bibitem{BO98}
%J.~Bourgain, 
%{\it Refinements of Strichartz' inequality and applications to 2D-NLS with critical nonlinearity},
% Internat. Math. Res. Notices 1998, no. 5, 253--283. 
%



\bibitem{Bring}
B.~Bringmann, 
{\it Invariant Gibbs measures for the three-dimensional wave equation with a Hartree nonlinearity II: dynamics},
 J. Eur. Math. Soc. (JEMS) 26 (2024), no. 6, 1933--2089. 







\bibitem{BDNY}
B.~Bringmann, Y.~Deng, A.R.~Nahmod, H.~Yue, 
{\it Invariant Gibbs measures for the three dimensional cubic nonlinear wave equation},
 Invent. Math. 236 (2024), no. 3, 1133--1411. 


\bibitem{BSch}
Y.~Bruned, K.~Schratz, 
{\it Resonance-based schemes for dispersive equations via decorated trees},
 Forum Math. Pi 10 (2022), Paper No. e2, 76 pp.





\bibitem{BT08}
N.~Burq, N.~Tzvetkov, 
{\it Random data Cauchy theory for supercritical wave equations. I. Local theory},
 Invent. Math. 173 (2008), no. 3, 449--475. 


\bibitem{BT3}
N.~Burq, N.~Tzvetkov, 
{\it Probabilistic well-posedness for the cubic wave equation},
 J. Eur. Math. Soc. (JEMS) 16 (2014), no. 1, 1--30. 
 
 





\bibitem{CatellierGubinelli}
R.~Catellier, M.~Gubinelli, 
{\it Averaging along irregular curves and regularisation of ODEs}, 
Stochastic Process. Appl. 126 (2016), no. 8, 2323--2366. 






\bibitem{CFLOP}
A.~Chapouto, J.~Forlano, G.~Li, T.~Oh, D.~Pilod,
{\it  Intermediate long wave equation in negative Sobolev spaces},
 Proc. Amer. Math. Soc. Ser. B 11 (2024), 452--468.


\bibitem{CFO}
A.~Chapouto, J.~Forlano, T.~Oh, 
{\it  Normal form integrator for the Korteweg-de Vries equation},
in preparation.


\bibitem{CGLLO2}
A.~Chapouto, M.~Gubinelli, G.~Li, J.~Li, T.~Oh, 
{\it Nonlinear PDEs with modulated dispersion III: multiplicative noises}, 
preprint.


\bibitem{CLOP}
A.~Chapouto, G.~Li, T.~Oh, D.~Pilod,
{\it Deep-water limit of the intermediate long wave equation in $L^2$},
Math. Res. Lett.
 31 (2024), no. 6, 1655--1692. 



\bibitem{CLO}
A.~Chapouto, G.~Li, T.~Oh, 
{\it
Deep-water and shallow-water limits
of statistical equilibria
for the intermediate long wave equation}, 
arXiv:2409.06905 [math.AP].




\bibitem{CLO2}
A.~Chapouto, J.~Li, T.~Oh, 
{\it  Fourier restriction norm method adapted to controlled paths:
stochastic wave equations},
preprint.






\bibitem{COZ}
A.~Chapouto, T.~Oh, G.~Zheng,
{\it  Pathwise well-posedness of the stochastic nonlinear Schr\"odinger equations with multiplicative noises},
in preparation.


%\bibitem{CGess}
%K.~Chouk, B.~Gess, 
%{\it Path-by-path regularization by noise for scalar conservation laws}, 
%J. Funct. Anal. 277 (2019), no. 5, 1469--1498. 
%



\bibitem{CG1}
K.~Chouk, M.~Gubinelli, 
{\it Nonlinear PDEs with modulated dispersion I: Nonlinear Schr\"odinger equations},
Comm. Partial Differential Equations 40 (2015), no. 11, 2047--2081. 





\bibitem{CGLLO1}
K.~Chouk, M.~Gubinelli, G.~Li, J.~Li, T.~Oh, 
{\it Nonlinear PDEs with modulated dispersion II: Korteweg-de Vries equation}, 
arXiv:1406.7675 [math.AP].




\bibitem{Christ}
M.~Christ, 
{\it Power series solution of a nonlinear Schr\"odinger equation}, Mathematical aspects of nonlinear dispersive equations, 131--155, Ann. of Math. Stud., 163, Princeton Univ. Press, Princeton, NJ, 2007.



\bibitem{CCT}
M.~Christ, J.~Colliander, T.~Tao,
{\it Asymptotics, frequency modulation, and low regularity ill-posedness for canonical defocusing equations},
Amer. J. Math. 125 (2003), no. 6, 1235--1293.





\bibitem{CGKO}
J.~Chung, Z.~Guo, S.~Kwon, T.~Oh, 
{\it  Normal form approach to global well-posedness of the quadratic derivative nonlinear Schr\"odinger equation on the circle}, Ann. Inst. H. Poincar\'e Anal. Non Lin\'eaire 34 (2017), 1273--1297. 





\bibitem{CKSTT03}
J.~Colliander, M.~Keel, G.~Staffilani, H.~Takaoka, T.~Tao,
{\it Sharp global well-posedness for KdV and modified KdV on $\R$ and $\T$},
J. Amer. Math. Soc. 16 (2003), no. 3, 705--749.




\bibitem{CKSTT4}
J.~Colliander, M.~Keel, G.~Staffilani, H.~Takaoka, T.~Tao,
{\it Resonant decompositions and the I-method for the cubic nonlinear Schr\"odinger equation on $\R^2$},  Discrete Contin. Dyn. Syst. 21 (2008), no. 3, 665--686. 

\bibitem{CO}
J.~Colliander, T.~Oh, 
{\it  Almost sure well-posedness of the cubic nonlinear Schr\"odinger equation below $L^2(\T)$}, 
Duke Math. J. 161 (2012), no. 3, 367--414. 



\bibitem{DD}
A.~de Bouard, A.~Debussche, 
{\it The nonlinear Schr\"odinger equation with white noise dispersion},
J. Funct. Anal. 259 (2010), no. 5, 1300--1321. 

%\bibitem{ddt-2} 
%A.~de Bouard, A.~Debussche, Y.~Tsutsumi,
%{\it White noise driven Korteweg-de Vries equation}, 
%J. Funct. Anal. 169 (1999), no. 2, 532--558



%\bibitem{DDT1} 
%A.~de Bouard, A.~Debussche, Y.~Tsutsumi,
%{\it Periodic solutions of the Korteweg-de Vries equation driven by white noise},
% SIAM J. Math. Anal. 36 (2004/05), no. 3, 815--855.
% 
% 
% 
% \bibitem{BLLZ}
%E.~Brun, G.~Li, R.~Liu, Y.~Zine,
%{\it  Global well-posedness of one-dimensional cubic fractional nonlinear Schr\"odinger equations in negative Sobolev spaces}, 
% arXiv:2311.13370 [math.AP].
%
 

\bibitem{DT}
A.~Debussche, Y.~Tsutsumi, 
{\it 1D quintic nonlinear Schr\"odinger equation with white noise dispersion},
 J. Math. Pures Appl.  96 (2011), no. 4, 363--376. 


\bibitem{Deng15}
Y.~Deng,
{\it Invariance of the Gibbs measure for the Benjamin-Ono equation}, 
J. Eur. Math. Soc. (JEMS) 17 (2015), no. 5, 1107--1198.



   \bibitem{DNY3}Y.~Deng, A.~Nahmod,  H.~Yue,
  {\it Random tensors, propagation of randomness, and
  nonlinear dispersive equations}, 
  Invent.
  Math. 228 (2022), no. 2, 539--686.



\bibitem{DNY2}
Y.~Deng, A.R.~Nahmod, H.~Yue,
{\it Invariant Gibbs measures and global strong solutions for nonlinear Schr\"odinger equations in dimension two},
Ann. of Math.  200 (2024), no. 2, 399--486.

 
%\bibitem{DGT12}
%A.~Deya, M.~Gubinelli, S.~Tindel,
%{\it Non-linear rough heat equations},
%Probab. Theory Related Fields 153 (2012), no. 1-2, 97--147.
 
 

\bibitem{DR2}
R.~Duboscq, A.~R\'eveillac, 
{\it On a stochastic Hardy-Littlewood-Sobolev inequality with application to Strichartz estimates for a noisy dispersion}, 
Ann. H. Lebesgue 5 (2022), 263--274. 


\bibitem{ETz}
M.B.~Erdo\u{g}an, N.~Tzirakis, 
{\it Global smoothing for the periodic KdV evolution}, 
Int. Math. Res. Not. IMRN (2013), no. 20, 4589--4614.


\bibitem{ETz3}
M.B.~Erdo\u{g}an, N.~Tzirakis, 
{\it Smoothing and global attractors for the Zakharov system on the torus},
 Anal. PDE 6 (2013), no. 3, 723--750.


\bibitem{ETz2}
M.B.~Erdo\u{g}an, N.~Tzirakis, 
{\it Dispersive partial differential equations. Wellposedness and applications.} London Mathematical Society Student Texts, 86. Cambridge University Press, Cambridge, 2016. xvi+186 pp.


\bibitem{Flan}
F.~Flandoli, 
{\it Random perturbation of PDEs and fluid dynamic models},
 Lectures from the 40th Probability Summer School held in Saint-Flour, 2010. Lecture Notes in Math., 2015. 
 \'Ec. \'Et\'e Probab. St.-Flour [Saint-Flour Probability Summer School] Springer, Heidelberg, 2011. x+176 pp.




\bibitem{FGP}
F.~Flandoli, M.~Gubinelli, E.~Priola,
{\it Well-posedness of the transport equation by stochastic perturbation},
Invent. Math. 180 (2010), no. 1, 1--53.





\bibitem{FLZ}
J.~Forlano, G.~Li, T.~Zhao,
{\it Unconditional deep-water limit of the intermediate long wave equation in low-regularity},
 NoDEA Nonlinear Differential Equations Appl. 32 (2025), no. 2, Paper No. 28, 31 pp.


\bibitem{FO}
J.~Forlano, T.~Oh, 
{\it  Normal form approach to the one-dimensional cubic nonlinear Schr\"odinger equation in Fourier-amalgam spaces}, preprint.



\bibitem{FOW}
J.~Forlano, T.~Oh, Y.~Wang,
{\it  Stochastic cubic nonlinear Schr\"odinger equation with almost space-time white noise}, J. Aust. Math. Soc. 109 (2020), no. 1, 44--67. 


\bibitem{FH20}
P.K.~Friz,  M.~Hairer, 
{\it A course on rough paths}, 
With an introduction to regularity structures. Second edition.
Universitext, Springer, Cham, [2020] \copyright 2020. xvi+346 pp.



%
%\bibitem{FrizVictoir}
%P.K.~Friz, N.B.~Victoir,
%{\it Multidimensional stochastic processes as rough paths},
%Theory and applications.
%Cambridge Stud. Adv. Math., 120.
%Cambridge University Press, Cambridge, 2010. xiv+656 pp.
%

\bibitem{G22}
L.~Galeati, 
{\it Pathwise methods in regularisation by noise}. - Bonn, 2022. - Dissertation, 
Rheinische Friedrich-Wilhelms-Universit\"at Bonn. 

\bibitem{G23}
L.~Galeati, 
{\it Nonlinear {Y}oung differential equations: a review}, 
J. Dynam. Differential Equations 35 (2023), no. 2, 985--1046.



\bibitem{GG0}
L.~Galeati, M.~Gubinelli,
{\it Noiseless regularisation by noise}, 
Rev. Mat. Iberoam.38(2022), no.2, 433--502.




\bibitem{GG}
L.~Galeati, M.~Gubinelli,
{\it  Prevalence of $\rho$-irregularity and related properties},
 Ann. Inst. H. Poincar\'e Probab. Stat.
 60 (2024), no. 4, 2415--2467.



%
%
%\bibitem{geman_occupation_1980}
%D.~Geman, J.~Horowitz,
%{\it Occupation densities},
%Ann. Probab. 8 (1980), no. 1, 1--67.
%



\bibitem{GKT23}
P.~G\'erard, T.~Kappeler, P.~Topalov,
{\it  Sharp well-posedness results of the Benjamin-Ono equation in $H^s(\T, \R)$  and qualitative properties of its solutions}, 
Acta Math. 231 (2023), no. 1,  31--88.



\bibitem{GKT24}
P.~G\'erard, T.~Kappeler, P.~Topalov,
{\it On smoothing properties and Tao's gauge transform of the Benjamin-Ono equation on the torus},
 Ann. Sci. \'Ec. Norm. Sup\'er.  57 (2024), no. 4, 1233--1270.



\bibitem{Gess}
B.~Gess, 
{\it Regularization and well-posedness by noise for ordinary and partial differential equations},
 Stochastic partial differential equations and related fields, 43--67. Springer Proc. Math. Stat., 229
Springer, Cham, 2018.




\bibitem{GTV}
J.~Ginibre, Y.~Tsutsumi, G.~Velo, 
{\it On the Cauchy problem for the Zakharov system}, 
J. Funct. Anal. 151 (1997), no. 2, 384--436.

%
%\bibitem{Grim}
%R.~Grimshaw, 
%{\it Internal Solitary Waves},
%Environmental Stratified Flows, 1--27,  Topics in Environmental Fluid Mechanics, vol 3. 
%(2003). 
%Springer, Boston, MA. 
%
%{\it Internal Solitary Waves}, Advances in Coastal and Ocean Engineering, ed. P.L.-F. Liu, World Scientific Publishing Company, Singapore, 3, 1 (1997).

%\bibitem{Gromov}
%M.~Gromov, 
%{\it Pseudo holomorphic curves in symplectic manifolds},
%Invent. Math. 82 (1985), no. 2, 307--347.


\bibitem{GH}
A.~Gr\"unrock, S.~Herr, 
{\it Low regularity local well-posedness of the derivative nonlinear Schr\"odinger equation with periodic initial data}, SIAM J. Math. Anal. 39 (2008), no. 6, 1890--1920.


\bibitem{Gub04}
M.~Gubinelli,
{\it Controlling rough paths},
J. Funct. Anal. 216 (2004), no. 1, 86--140.



\bibitem{GubinelliKdV}
M.~Gubinelli, 
{\it Rough solutions for the periodic Korteweg-de Vries equation},
Commun. Pure Appl. Anal. 11 (2012), no. 2, 709--733.


%\bibitem{GLT}
%M.~Gubinelli, A.~Lejay, S.~Tindel,
%{\it Young integrals and SPDEs},
%Potential Anal. 25 (2006), no. 4, 307--326.



%\bibitem{GIP}
%M.~Gubinelli, P.~Imkeller, N.~Perkowski,
%{\it Paracontrolled distributions and singular PDEs},
%Forum Math. Pi 3 (2015), e6, 75 pp.

%
%
%\bibitem{gubinelli_regularization_2012}
%M.~Gubinelli, M.~Jara,
%{\it Regularization by noise and stochastic Burgers equations},
%Stoch. Partial Differ. Equ. Anal. Comput.1 (2013), no. 2, 325--350.
%

\bibitem{GKO2}
M.~Gubinelli, H.~Koch, T.~Oh, 
{\it  Paracontrolled approach to the three-dimensional stochastic nonlinear wave equation with quadratic nonlinearity},  
J. Eur. Math. Soc. (JEMS) 26 (2024), no. 3, 817--874. 



\bibitem{GT10}
M.~Gubinelli, S.~Tindel,
{\it Rough evolution equations},
Ann. Probab. 38 (2010), no. 1, 1--75.




\bibitem{GKO}
Z.~Guo, S.~Kwon, T.~Oh, 
{\it Poincar\'e-Dulac normal form reduction for unconditional well-posedness of
the periodic cubic NLS}, Comm. Math. Phys. 322 (2013), no. 1, 19--48.


\bibitem{GO}
Z.~Guo,  T.~Oh, 
{\it Non-existence of solutions for the periodic cubic nonlinear Schr\"odinger equation below $L^2$}, Int. Math. Res. Not. 2018, no. 6, 1656--1729. 




\bibitem{GP}
I.~Gy\"ongy,  \'E.~Pardoux, 
    {\it On the regularization effect of space-time white noise on quasi-linear parabolic partial differential equations},
    Probab. Theory Related Fields 97 (1993), no. 1--2, 211--229.



%\bibitem{HW}
%G.H.~Hardy, E.M.~Wright, {\it An introduction to the theory of numbers,} Fifth edition.
%	The Clarendon Press, Oxford University Press, New York, 1979. xvi+426 pp. 


\bibitem{HMSch}
M.~Hofmanov\'a, M.~Kn\"oller, K.~Schratz,
{\it Randomized exponential integrators for modulated nonlinear Schr\"odinger equations}, IMA J. Numer. Anal. 40 (2020), no. 4, 2143--2162

\bibitem{HSch}
M.~Hofmanov\'a,  K.~Schratz,
{\it An exponential-type integrator for the KdV equation},  Numer. Math. 136 (2017), no. 4, 1117--1137.


\bibitem{HK}
Y.~Hu, K.~L\^e, 
{\it Nonlinear Young integrals and differential systems in H\"older media},
 Trans. Amer. Math. Soc. 369 (2017), no. 3, 1935--2002. 


\bibitem{KT1}
T.~Kappeler, P.~Topalov,
{\it Global wellposedness of KdV in $H^{-1}(\T, \R)$},
Duke Math. J. 135 (2006), no. 2, 327--360.




\bibitem{KATO}
T.~Kato, {\it On nonlinear Schr\"odinger equations. II. $H^s$-solutions and unconditional
well-posedness,} J. Anal. Math. 67 (1995), 281--306.




%\bibitem{KPV91}
%C.E.~Kenig, G.~Ponce, L.~Vega,
%{\it Well-posedness of the initial value problem for the Korteweg-de Vries equation},
%J. Amer. Math. Soc. 4 (1991), no. 2, 323--347.
%
%
%
%\bibitem{KPV93}
%C.E.~Kenig, G.~Ponce, L.~Vega,
%{\it Well-posedness and scattering results for the generalized Korteweg-de Vries equation via the contraction principle},
%Comm. Pure Appl. Math.46 (1993), no.4, 527--620.




\bibitem{KPV96}
C.E.~Kenig, G.~Ponce, L.~Vega,
{\it A bilinear estimate with applications to the KdV equation},
J. Amer. Math. Soc. 9 (1996), no. 2, 573--603.






\bibitem{KLV}
R.~Killip, T.~Laurens, M.~Vi\c{s}an, 
{\it Sharp well-posedness for the Benjamin--Ono equation}, 
Invent. Math. 236 (2024), no. 3, 999--1054.





\bibitem{KV19}
R.~Killip, M.~Vi\c{s}an, 
{\it KdV is well-posed in $H^{-1}$}, 
Ann. of Math. 190 (2019), no.1, 249--305.



\bibitem{Ki2}
N.~Kishimoto,
{\it Unconditional uniqueness of solutions for nonlinear dispersive equations}, 
arXiv:1911.04349 [math.AP].


\bibitem{Ki3}
N.~Kishimoto,
{\it Unconditional uniqueness for the periodic Benjamin-Ono equation by normal form approach},
 J. Math. Anal. Appl. 514 (2022), no. 2, Paper No. 126309, 22 pp.


\bibitem{KR}
N.V.~Krylov, M.~R\"ockner,
    {\it Strong solutions of stochastic equations with singular time dependent drift},
    Probab. Theory Related Fields 131 (2005), no. 2, 154--196. 





\bibitem{KO}
S.~Kwon, T.~Oh, 
{\it  On unconditional well-posedness of modified KdV}, Int. Math. Res. Not. (2012), no. 15, 3509--3534. 






\bibitem{KOY}
S.~Kwon, T.~Oh, H.~Yoon, 
{\it Normal form approach to unconditional well-posedness of nonlinear dispersive PDEs on the real line}, Ann. Fac. Sci. Toulouse Math. 29 (2020), no. 3, 649--720. 



\bibitem{Li24}
G.~Li, 
{\it Deep-water and shallow-water limits of the intermediate long wave equation}, 
 Nonlinearity 37 (2024), no.7, Paper No. 075001, 44 pp. 
 
 
\bibitem{LOZ}
G.~Li, T.~Oh, G.~Zheng,
{\it On the deep-water and shallow-water limits of the intermediate long wave equation from a statistical viewpoint}, 
Trans. London Math. Soc. 
 12 (2025), no. 1, Paper No. e70005. 






%
%
%\bibitem{marty2006}
%R.~Marty,
%{\it On a splitting scheme for the nonlinear Schr\"odinger equation in a random medium},
%Commun. Math. Sci. 4 (2006), no. 4, 679--705.
%


%
%
%\bibitem{Attila}
%P.~Mattila,
%{\it  Hausdorff dimension, projections, and the Fourier transform},
%Publ. Mat. 48 (2004), no. 1, 3--48.
%
%



%
%\bibitem{M68}
%R.M.~Miura,
%{\it  Korteweg-de Vries equation and generalizations. \textup{I}. A remarkable explicit nonlinear transformation},
%J. Mathematical Phys. 9 (1968), 1202--1204.
%



\bibitem{Moli1}
L.~Molinet, 
{\it Global well-posedness in the energy space for the Benjamin-Ono equation on the circle},
Math. Ann. 337 (2007), no. 2, 353--383. 


\bibitem{Moli2}
L.~Molinet, 
{\it Global well-posedness in $L^2$ for the periodic Benjamin-Ono equation}, 
Amer. J. Math. 130 (2008), no. 3, 635--683. 








\bibitem{M12}
L.~Molinet,
{\it Sharp ill-posedness results for the KdV and mKdV equations on the torus},
Adv. Math. 230 (2012), no. 4--6, 1895--1930.




%\bibitem{MP}
%L.~Molinet, D.~Pilod, 
%{\it The Cauchy problem for the Benjamin-Ono equation in $L^2$ revisited},
% Anal. PDE 5 (2012), no. 2, 365--395. 



\bibitem{MST}
L.~Molinet, J.-C.~Saut, N.~Tzvetkov, 
{\it Ill-posedness issues for the Benjamin-Ono and related equations},
SIAM J. Math. Anal. 33 (2001), no. 4, 982--988. 


\bibitem{MP2}
R.~Mo\c{s}incat, D.~Pilod, 
{\it  Unconditional uniqueness for the Benjamin-Ono equation}, Pure Appl. Anal. 5 (2023), no. 2, 285--322. 



%
%\bibitem{NTT}
%K.~Nakanishi, H.~Takaoka, Y.~Tsutsumi, 
%{\it Local well-posedness in low regularity of the mKdV equation with periodic boundary condition},
% Discrete Contin. Dyn. Syst. 28 (2010), no. 4, 1635--1654. 
 

\bibitem{OOT1}
T.~Oh, M.~Okamoto, L.~Tolomeo,
{\it Focusing $\Phi^4_3$-model with a Hartree-type nonlinearity}, 
Mem. Amer. Math. Soc. 304 (2024), no. 1529, vi+143 pp. 



\bibitem{OOT2}
T.~Oh, M.~Okamoto, L.~Tolomeo,
{\it Stochastic quantization of the $\Phi^3_3$-model}, 
Mem. Eur. Math. Soc., 16. EMS Press, Berlin, 2025, viii+145 pp. 





\bibitem{OP}
T.~Oh, O.~Pocovnicu,
{\it  Probabilistic global well-posedness of the energy-critical defocusing quintic nonlinear wave equation on $\R^d$}, J. Math. Pures Appl. 105 (2016), 342--366. 



\bibitem{OS}
T.~Oh, K.~Seong,
{\it  Quasi-invariant Gaussian measures for the cubic fourth order nonlinear Schr\"odinger equation in negative Sobolev spaces}, J. Funct. Anal. 281 (2021), no. 9, 109150, 49 pp. 


\bibitem{OST}
T.~Oh, P.~Sosoe, N.~Tzvetkov,
{\it  An optimal regularity result on the quasi-invariant Gaussian measures for the cubic fourth order nonlinear Schr\"odinger equation},  
J. \'Ec. polytech. Math. 5 (2018), 793--841. 

\bibitem{OSW}
T.~Oh, P.~Sosoe, Y.~Wang, 
{\it Normal form approach to the stochastic Korteweg-de Vries
equation with an additive noise}, 
in prepration.




\bibitem{OW2}
T.~Oh, Y.~Wang,
{\it Global well-posedness of the periodic cubic fourth order NLS in negative Sobolev spaces},
 Forum Math. Sigma 6 (2018), e5, 80 pp. 



\bibitem{OW1}
T.~Oh, Y.~Wang,
{\it Global well-posedness of the one-dimensional cubic nonlinear Schr\"odinger equation in almost critical spaces}, J. Differential Equations 269 (2020), no. 1, 612--640. 


\bibitem{OW}
T.~Oh, Y.~Wang,
{\it  Normal form approach to the one-dimensional periodic cubic nonlinear Schr\"odinger equation in almost critical Fourier--Lebesgue spaces}, 
J. Anal. Math. 143 (2021), no. 2, 723--762. 

\bibitem{OW3}
T.~Oh, Y.~Wang,
{\it Revisiting Bourgain's probabilistic construction of solutions to the 2-d cubic NLS}, 
preprint.


\bibitem{OWZ}
T.~Oh, Y.~Wang, Y.~Zine,
\textit{Three-dimensional stochastic cubic nonlinear wave equation with almost space-time white noise},
Stoch. Partial Differ. Equ. Anal. Comput. 10 (2022), no. 3, 898--963.


\bibitem{Poc}
O.~Pocovnicu, 
{\it Almost sure global well-posedness for the energy-critical defocusing nonlinear wave equation on $\R^d$,
 $d = 4$ and $5$}, J. Eur. Math. Soc. (JEMS) 19 (2017), no. 8, 2521--2575.






\bibitem{Robert1}
T.~Robert, 
{\it Regularization by noise for some nonlinear dispersive PDEs},
 Journ\'ees \'equations aux d\'eriv\'ees partielles (2023), Talk no. 9, 12 p. 





\bibitem{Robert2}
T.~Robert, 
{\it Remarks on nonlinear dispersive PDEs with rough dispersion management}, 
arXiv:2410.23038 [math.AP].




\bibitem{Robert3}
T.~Robert, 
{\it Regularization by noise for some strongly non-resonant modulated dispersive PDEs}, 
arXiv:2410.23051 [math.AP].


\bibitem{Robert4}
T.~Robert, 
{\it  On the Kadomtsev-Petviashvili equations with rough dispersion management}, 
 in preparation.



\bibitem{RT}
M.~Romito, L.~Tolomeo, 
{\it Yet another notion of irregularity through small ball estimates},
arXiv:2207.02716 [math.PR].


	
\bibitem{S19}
J.-C. Saut,
 {\it Benjamin-Ono and intermediate long wave equations: modeling, IST and PDE. Nonlinear
dispersive partial differential equations and inverse scattering}, 
95--160, Fields Inst. Commun., 83, Springer, New York, [2019], \textcopyright2019.


\bibitem{Ste}
G.~Stewart, 
{\it On the wellposedness of periodic nonlinear Schr\"odinger equations with white noise dispersion}, 
Stoch. Partial Differ. Equ. Anal. Comput. 12 (2024), no. 3, 1417--1438.


\bibitem{TT}
H.~Takaoka, Y.~Tsutsumi, 
{\it Well-posedness of the Cauchy problem for the modified KdV equation with periodic boundary condition},
 Int. Math. Res. Not. (2004), no. 56, 3009--3040. 




\bibitem{Tanaka}
T.~Tanaka, 
{\it Remark on the local well-posedness for NLS with the modulated dispersion}, 
arXiv:2407.20005 [math.AP].


\bibitem{Tao2}
T.~Tao, 
{Multilinear weighted convolution of $L^2$-functions, and applications to non-
linear dispersive equations}, Amer. J. Math. 123 (2001), no. 5, 839--908.








\end{thebibliography}







\end{document}






